% Generated mechanically from frozen input inputs/p09/ja/src/spaces-over-fields.tex.
% Do not edit this generated file; edit the bound locale source instead.

% OK, start here.
%

\setcounter{chapter}{71}
\chapter{体上の代数空間}



\phantomsection
\label{spaces-over-fields-section-phantom}


\section{はじめに}
\label{spaces-over-fields-section-introduction}

\noindent
本章は、多様体に関する章を代数空間の枠組みに移したものである。
代数空間についての参考文献として \cite{Kn} がある。


\section{約束}
\label{spaces-over-fields-section-conventions}

\noindent
すべてのスキームは大 fppf サイト $\Sch_{fppf}$ に含まれるものとする。
また、考えるすべての環 $A$ は、$\Spec(A)$ がこの大サイトの対象と
（同型で）あるという性質をもつものとする。

\medskip\noindent
$S$ をスキームとし、$X$ を $S$ 上の代数空間とする。
本章および後続の章では $X \times_S X$ と書き、これは $X$ とそれ自身との
（$S$ 上の代数空間の圏における）積を表す。$X \times X$ とは書かない。



\section{一般有限射}
\label{spaces-over-fields-section-generically-finite}

\noindent
本節では、Decent Spaces, Section
\ref{decent-spaces-section-generically-finite} の議論と、
Varieties, Section \ref{varieties-section-generically-finite}
における代数空間の射についての対応物を続ける。

\begin{lemma}
\label{spaces-over-fields-lemma-quasi-finite-in-codim-1}
$S$ をスキームとする。$f : X \to Y$ を $S$ 上の代数空間の射とする。
$f$ は局所有限型であり、$Y$ は局所 Noether であると仮定する。
$y \in |Y|$ を余次元 $\leq 1$ の $Y$ の点とする。
$X^0 \subset |X|$ を余次元 $0$ の $X$ の点全体とする。
さらに、次の条件のいずれかが成り立つと仮定する。
\begin{enumerate}
\item すべての $x \in X^0$ に対して $x/f(x)$ の超越次数は $0$ である。
\item すべての $x \in X^0$ で $f(x) \leadsto y$ を満たすものに対して、
$x/f(x)$ の超越次数は $0$ である。
\item $f$ はすべての $x \in X^0$ において準有限である。
\item $f$ は $|X|$ の稠密な点集合において準有限である。
% P09-ERR-0241
\item ここにさらに追加する。
\end{enumerate}
このとき、$f$ は $X$ の $y$ 上のすべての点において準有限である。
\end{lemma}

\begin{proof}
証明をスキームの場合に帰着させる。そのため、可換図式
$$
\xymatrix{
U \ar[r] \ar[d]_g & X \ar[d]^f \\
V \ar[r] & Y
}
$$
を選ぶ。ここで $U$, $V$ はスキームであり、横向きの射は \'etale かつ全射である。
$v \in V$ を $y$ に写る点として選ぶ。$V$ は局所 Noether であり、
$\dim(\mathcal{O}_{V, v}) \leq 1$ であることに注意せよ
（Properties of Spaces, Definitions
\ref{spaces-properties-definition-dimension-local-ring} および Remark
\ref{spaces-properties-remark-list-properties-local-etale-topology} を参照）。
$U_v$ は $U \to V$ の $v$ 上のファイバーであり、
$f^{-1}(\{y\}) \subset |X|$ へ全射である。$X^0$ の $U$ における逆像は、
$U$ の既約成分の一般点全体にちょうど一致する（Properties of Spaces, Lemma
\ref{spaces-properties-lemma-codimension-0-points}）。
$\eta \in U$ をそのような点とし、その像を $x \in X^0$ とすると、
$x / f(x)$ の超越次数は $\kappa(\eta)$ の $\kappa(g(\eta))$ 上の超越次数である
（Morphisms of Spaces, Definition
\ref{spaces-morphisms-definition-dimension-fibre}）。
$U \to V$ が $u \in U$ において準有限であることと、$f$ が $u$ の $X$ における像で
準有限であることは同値である。

\medskip\noindent
場合 (1)。この場合、Varieties, Lemma
\ref{varieties-lemma-quasi-finite-in-codim-1} の場合 (1) を適用でき、
$U \to V$ は $U_v$ のすべての点において準有限である。
従って、$f$ は $y$ の上にあるすべての点で準有限である。

\medskip\noindent
場合 (2)。$u \in U$ を、$V$ における像が $v$ に特殊化する既約成分の一般点とする。
すると、$x \in X^0$ を $u$ の像とすれば、$f(x) \leadsto y$ が成り立つ。
従って、Varieties, Lemma
\ref{varieties-lemma-quasi-finite-in-codim-1} の場合 (2) を適用でき、
前と同じ結論を得る。

\medskip\noindent
場合 (3) は Varieties, Lemma
\ref{varieties-lemma-quasi-finite-in-codim-1} の場合 (3) から従う。

\medskip\noindent
場合 (4) では、$|U| \to |X|$ は開写像なので、
$U \to V$ が準有限である点の集合も稠密である。
従って、Varieties, Lemma
\ref{varieties-lemma-quasi-finite-in-codim-1} の場合 (4) を適用できる。
\end{proof}

\begin{lemma}
\label{spaces-over-fields-lemma-finite-in-codim-1}
$S$ をスキームとする。$f : X \to Y$ を $S$ 上の代数空間の射とする。
$f$ は固有であり、$Y$ は局所 Noether であると仮定する。
$y \in Y$ を、余次元が $\leq 1$ である $Y$ の点とする。
$X^0 \subset |X|$ を余次元 $0$ の $X$ の点全体とする。
さらに、次の条件のいずれかが成り立つと仮定する。
\begin{enumerate}
\item すべての $x \in X^0$ に対して $x/f(x)$ の超越次数は $0$ である。
\item すべての $x \in X^0$ で $f(x) \leadsto y$ を満たすものに対して、
$x/f(x)$ の超越次数は $0$ である。
\item $f$ はすべての $x \in X^0$ において準有限である。
\item $f$ は $|X|$ の稠密な点集合において準有限である。
% P09-ERR-0241
\item ここにさらに追加する。
\end{enumerate}
このとき、開部分空間 $Y' \subset Y$ で $y$ を含み、
$Y' \times_Y X \to Y'$ が有限となるものが存在する。
\end{lemma}

\begin{proof}
Lemma \ref{spaces-over-fields-lemma-quasi-finite-in-codim-1} により、射 $f$ は $y$ の上にある
すべての点において準有限である。$\overline{y} : \Spec(k) \to Y$ を
$y$ の上にある幾何学的点とする。このとき $|X_{\overline{y}}|$ は離散空間である
（Decent Spaces, Lemma
\ref{decent-spaces-lemma-conditions-on-fibre-and-qf}）。
$X_{\overline{y}}$ は、$f$ が固有なので準コンパクトであり、従って
$|X_{\overline{y}}|$ は有限である。ゆえに Cohomology of Spaces, Lemma
\ref{spaces-cohomology-lemma-proper-finite-fibre-finite-in-neighbourhood}
を適用して結論を得る。
\end{proof}

\begin{lemma}
\label{spaces-over-fields-lemma-modification-normal-iso-over-codimension-1}
$S$ をスキームとする。$X$ を $S$ 上の Noether 代数空間とする。
$f : Y \to X$ を代数空間の双有理固有射とし、$Y$ は被約であるとする。
$U \subset X$ を、その上で $f$ が同型となる最大の開部分空間とする。
このとき $U$ は次を含む。
\begin{enumerate}
\item 余次元 $0$ の $X$ のすべての点。
\item $x \in |X|$ であって、余次元が $1$（$X$ 上）であり、$X$ の $x$ における局所環が
正規であるもの（Properties of Spaces, Remark
\ref{spaces-properties-remark-list-properties-local-ring-local-etale-topology}）。
\item すべての $x \in |X|$ で、$|Y| \to |X|$ の $x$ 上のファイバーが有限であり、
かつ $X$ の $x$ における局所環が正規であるもの。
\end{enumerate}
\end{lemma}

\begin{proof}
(1) は Decent Spaces, Lemma
\ref{decent-spaces-lemma-birational-isomorphism-over-dense-open}
から従う（Noether 代数空間 $X$ と $Y$ は準分離であり、従って decent であることも用いる）。
(2) は (3) と Lemma \ref{spaces-over-fields-lemma-finite-in-codim-1} から従う
（有限射のファイバーが有限であることも用いる）。
(3) のような $x \in |X|$ をとる。Cohomology of Spaces, Lemma
\ref{spaces-cohomology-lemma-proper-finite-fibre-finite-in-neighbourhood}
（Decent Spaces, Lemma
\ref{decent-spaces-lemma-conditions-on-fibre-and-qf} により適用できる）から、
$f$ は有限であると仮定してよい。アフィンスキーム $X'$ と \'etale 射
$X' \to X$、および
% P09-ERR-0242
点 $x' \in X$ を選び、その像が $x$ であるとする。
$U'$ という $x' \in X'$ の開近傍で、$Y \times_X X' \to X'$ が $U'$ 上で同型となるものが
存在することを示せば十分である（そのとき $U$ は $U'$ の $X$ における像を含む。
Spaces, Lemma
\ref{spaces-lemma-descent-representable-transformations-property} を参照）。
すると
% P09-ERR-0243
$Y \times_X X' \to X$ は有限双有理射である
（Decent Spaces, Lemma
\ref{decent-spaces-lemma-birational-etale-localization}）。
有限射はアフィンなので、Noether アフィンスキームの有限双有理射
$Y \to X$ と、$x \in X$ で $\mathcal{O}_{X, x}$ が正規整域となるものの場合に帰着する。
これは Varieties, Lemma
\ref{varieties-lemma-modification-normal-iso-over-codimension-1}
で扱われている。
\end{proof}






\section{整代数空間}
\label{spaces-over-fields-section-integral-spaces}

\noindent
整代数空間という概念はまだ定義していない。問題は、整であることがスキームの
\'etale 局所的性質ではないことである。Properties, Lemma
\ref{properties-lemma-characterize-integral} にある、$X$ が被約で $|X|$ が既約である
という性質を用いて整代数空間を定義することもできる。しかし、その場合には Spaces,
Example \ref{spaces-example-infinite-product} で述べた代数空間 $X$ が整となり、
これは適切とは思われない。この種の病理を避けるため、より弱い代替条件があるかもしれないが、
さらにそれが decent な代数空間であると仮定する。

\begin{definition}
\label{spaces-over-fields-definition-integral-algebraic-space}
$S$ をスキームとする。代数空間 $X$ が $S$ 上で、被約かつ decent であり、
$|X|$ が既約であるとき、これは {\it 整} であるという。
\end{definition}

\noindent
この場合、既約位相空間 $|X|$ は sober である
（Decent Spaces, Proposition
\ref{decent-spaces-proposition-reasonable-sober}）。
従って、この空間は一意な一般点 $x$ をもつ。実際、Decent Spaces, Lemma
\ref{decent-spaces-lemma-finitely-many-irreducible-components}
では、既約成分を有限個しかもたない decent 代数空間を特徴づけた。
この補題を適用すると、代数空間 $X$ が整であるための一つの条件は、被約であり、
既約稠密開部分スキーム $X'$ が一般点 $x'$ をもち、
射 $x' \to X$ が準コンパクトであることである。

\begin{lemma}
\label{spaces-over-fields-lemma-integral-algebraic-space-rational-functions}
$S$ をスキームとする。$X$ を $S$ 上の整代数空間とする。
$\eta \in |X|$ を $X$ の一般点とする。標準的な同一視
$$
R(X) = \mathcal{O}_{X, \eta}^h = \kappa(\eta)
$$
が存在する。ここで $R(X)$ は Morphisms of Spaces, Definition
\ref{spaces-morphisms-definition-ring-of-rational-functions}
で定義された有理関数環、$\kappa(\eta)$ は Decent Spaces, Definition
\ref{decent-spaces-definition-residue-field} で定義された剰余体、
$\mathcal{O}_{X, \eta}^h$ は Decent Spaces, Definition
\ref{decent-spaces-definition-elemenary-etale-neighbourhood}
で定義された Hensel 局所環である。特に、これらの環は体である。
\end{lemma}

\begin{proof}
$X$ は $\eta$ のある開近傍上でスキームである（上の議論を参照）ので、
これはスキームに対する対応結果から直ちに従う。Morphisms, Lemma
\ref{morphisms-lemma-integral-scheme-rational-functions} を参照せよ。
さらに、体の Hensel 化はその体自身であること、および代数空間に対するこれらの対象の
定義がスキームに対する定義と整合することを用いる。詳細は省略する。
\end{proof}

\noindent
これにより、次の定義を得る。

\begin{definition}
\label{spaces-over-fields-definition-function-field}
$S$ をスキームとする。$X$ を $S$ 上の整代数空間とする。
$X$ の {\it 関数体}、または {\it 有理関数体} とは、
Lemma \ref{spaces-over-fields-lemma-integral-algebraic-space-rational-functions} の体 $R(X)$ のことである。
\end{definition}

\noindent
この体を $k(X)$ と記し、$R(X)$ の代わりに用いることもある。

\begin{lemma}
\label{spaces-over-fields-lemma-integral-sections}
$S$ をスキームとする。$X$ を $S$ 上の整代数空間とする。
このとき $\Gamma(X, \mathcal{O}_X)$ は整域である。
\end{lemma}

\begin{proof}
$R = \Gamma(X, \mathcal{O}_X)$ とおく。$f, g \in R$ が非零で $fg = 0$ ならば、
$X = V(f) \cup V(g)$ である。ここで $V(f)$ は $X$ の、$f$ によって切り出される
閉部分空間を表す。$X$ は既約なので、$V(f) = X$ または $V(g) = X$ である。
従って Properties of Spaces, Lemma
\ref{spaces-properties-lemma-reduced-space} により、$f = 0$ または $g = 0$ である。
\end{proof}

\noindent
正規整代数空間に関する補題を述べる。

\begin{lemma}
\label{spaces-over-fields-lemma-normal-integral-cover-by-affines}
$S$ をスキームとする。$X$ を $S$ 上の正規整代数空間とする。
すべての $x \in |X|$ に対して、正規整アフィンスキーム $U$ と \'etale 射
$U \to X$ で、その像が $x$ を含むものが存在する。
\end{lemma}

\begin{proof}
アフィンスキーム $U$ と \'etale 射 $U \to X$ で、その像が $x$ を含むものを選ぶ。
% P09-ERR-0244
$u_i$, $i \in I$ を $U$ の既約成分の一般点とする。各 $u_i$ は $X$ の一般点へ写る
（Decent Spaces, Lemma
\ref{decent-spaces-lemma-decent-generic-points}）。
decent 空間の定義（Decent Spaces, Definition
\ref{decent-spaces-definition-very-reasonable}）により、$I$ は有限である。
従って $U = \Spec(A)$ であり、$A$ は極小素イデアルを有限個しかもたない正規環である。
Algebra, Lemma \ref{algebra-lemma-characterize-reduced-ring-normal} により、
$A = \prod_{i \in I} A_i$ は正規整域の積である。
従って $U = \coprod U_i$、$U_i = \Spec(A_i)$ であり、$x$ は
$U_i \to X$ の像に含まれる（ある $i$ に対して）。これで補題が証明された。
\end{proof}

\begin{lemma}
\label{spaces-over-fields-lemma-normal-integral-sections}
$S$ をスキームとする。$X$ を $S$ 上の正規整代数空間とする。
このとき $\Gamma(X, \mathcal{O}_X)$ は正規整域である。
\end{lemma}

\begin{proof}
$R = \Gamma(X, \mathcal{O}_X)$ とおく。Lemma
\ref{spaces-over-fields-lemma-integral-sections} により $R$ は整域である。
$f = a/b$ を、$R$ 上整な $R$ の分数体の元とする。
任意の \'etale 射 $U \to X$ で $U$ がスキームであるものに対して、
高々一つの $f_U \in \Gamma(U, \mathcal{O}_U)$ があり、これは $b|_U f_U = a|_U$ を満たす。
実際、$U$ は被約であり、$U$ の一般点は $X$ の一般点へ写るので、
$b|_U$ は零因子でない。すべての $x \in |X|$ に対し、Lemma
\ref{spaces-over-fields-lemma-normal-integral-cover-by-affines} のように $U \to X$ を選ぶ。
一意な $f_U \in \Gamma(U, \mathcal{O}_U)$ が存在し、これは
$b|_U f_U = a|_U$ を満たす。実際、$\Gamma(U, \mathcal{O}_U)$ は
Properties, Lemma \ref{properties-lemma-normal-integral-sections} により正規整域である。
上の一意性により、これらの $f_U$ は貼り合わさり、構造層の大域切断 $f$、すなわち
$R$ の元を定める。
\end{proof}

\begin{lemma}
\label{spaces-over-fields-lemma-decent-irreducible-closed}
$S$ をスキームとする。$X$ を $S$ 上の decent 代数空間とする。
次の集合の間には標準的な全単射がある。
\begin{enumerate}
\item $X$ の点の集合、すなわち $|X|$。
\item $|X|$ の既約閉部分集合の集合。
\item $X$ の整な閉部分空間の集合。
\end{enumerate}
(1) から (2) への全単射は $x$ を $\overline{\{x\}}$ へ写す。
(3) から (2) への全単射は $Z$ を $|Z|$ へ写す。
\end{lemma}

\begin{proof}
$|X|$ は Decent Spaces, Proposition
\ref{decent-spaces-proposition-reasonable-sober} により sober なので、
この写像は (1) と (2) の間の全単射を定める。
閉かつ既約な $T \subset |X|$ が与えられると、一意な被約閉部分空間
$Z \subset X$ で $|Z| = T$ を満たすものが存在する。すなわち、$Z$ は $T$ 上の被約誘導部分空間構造である。
Properties of Spaces, Definition
\ref{spaces-properties-definition-reduced-induced-space} を参照せよ。
これは decent、被約、かつ既約なので、整代数空間である。
\end{proof}







\section{整代数空間の間の射}
\label{spaces-over-fields-section-morphisms-between-integral-spaces}

\noindent
次の補題は、整代数空間の間の有限次数の支配射を特徴づける。

\begin{lemma}
\label{spaces-over-fields-lemma-finite-degree}
$S$ をスキームとする。
% P09-ERR-0245
$X$, $Y$ を $S$ 上の整代数空間とする。$x \in |X|$ と $y \in |Y|$ を一般点とする。
$f : X \to Y$ は局所有限型であるとする。$f$ が支配的であると仮定する
（Morphisms of Spaces, Definition \ref{spaces-morphisms-definition-dominant}）。
次の条件は同値である。
\begin{enumerate}
\item $x/y$ の超越次数が $0$ である。
\item 拡大 $\kappa(x)/\kappa(y)$（証明を参照）が有限である。
\item 非空アフィン開集合 $U \subset X$ と $V \subset Y$ が存在し、
$f(U) \subset V$ かつ $f|_U : U \to V$ が有限である。
\item $f$ が $x$ において準有限である。
\item $x$ が $|X|$ のうち $y$ に写る唯一の点である。
\end{enumerate}
$f$ が分離的であるか、または $f$ が準コンパクトであるならば、これらはさらに
\begin{enumerate}
\item[(6)] 非空アフィン開集合 $V \subset Y$ が存在し、$f^{-1}(V) \to V$ が有限である
\end{enumerate}
という条件とも同値である。
\end{lemma}

\begin{proof}
初等的な位相の議論により、$f(x) = y$ である。実際、$f$ は支配的である。
$Y' \subset Y$ を $Y$ のスキーム的軌跡とし、
$X' \subset f^{-1}(Y')$ を $f^{-1}(Y')$ のスキーム的軌跡とする。
上の議論に Decent Spaces, Proposition
\ref{decent-spaces-proposition-reasonable-sober} および Theorem
\ref{decent-spaces-theorem-decent-open-dense-scheme} を用いると、
$x \in |X'|$ かつ $y \in |Y'|$ であることが分かる。
このとき $f|_{X'} : X' \to Y'$ は局所有限型の整スキーム間の射である。
従って、Morphisms, Lemma \ref{morphisms-lemma-finite-degree} により
(1)、(2)、(3) は同値である。

\medskip\noindent
$X \to Y \to Y$ に Morphisms of Spaces, Lemma
\ref{spaces-morphisms-lemma-compare-tr-deg} を適用すれば、条件 (4) は条件 (1) を含意する。
一方、有限射は準有限であり、また $x \in U$ である（$x$ は一般点だからである）ので、
条件 (3) は条件 (4) を含意する。従って (1) -- (4) は同値である。

\medskip\noindent
同値な条件 (1) -- (4) を仮定する。$x' \mapsto y$ とする。このとき
$x \leadsto x'$ は $|X| \to |Y|$ の $y$ 上のファイバー内の特殊化である。
$x' \not = x$ ならば、Decent Spaces, Lemma
\ref{decent-spaces-lemma-conditions-on-point-in-fibre-and-qf} により、
$f$ は $x$ において準有限ではない。従って $x = x'$ であり、(5) が成り立つ。
逆に (5) が成り立つならば、スキームの射 $X' \to Y'$（上を参照）についても
(5) が成り立つので、Morphisms, Lemma \ref{morphisms-lemma-finite-degree}
を用いて (1) が成り立つことが分かる。

\medskip\noindent
(6) は、$f$ に何ら追加の仮定を課さなくても、同値な条件 (1) -- (5) を含意する。
証明を終えるには、同値な条件 (1) -- (5) が (6) を含意することを示せばよい。
これは Decent Spaces, Lemma
\ref{decent-spaces-lemma-finite-over-dense-open} から従う。
\end{proof}

\begin{definition}
\label{spaces-over-fields-definition-degree}
$S$ をスキームとする。$X$ と $Y$ を $S$ 上の整代数空間とする。
$f : X \to Y$ を局所有限型かつ支配的な射とする。Lemma
\ref{spaces-over-fields-lemma-finite-degree} の同値な条件 (1) -- (5) のいずれかを仮定する。
$x \in |X|$ と $y \in |Y|$ を一般点とする。このとき、正の整数
$$
\deg(X/Y) = [\kappa(x) : \kappa(y)]
$$
を {\it $X$ の $Y$ 上の次数} という。
\end{definition}

\begin{lemma}
\label{spaces-over-fields-lemma-degree-composition}
$S$ をスキームとする。$X$, $Y$, $Z$ を $S$ 上の整代数空間とする。
$f : X \to Y$ と $g : Y \to Z$ を局所有限型の支配射とする。Lemma
\ref{spaces-over-fields-lemma-finite-degree} の同値な条件 (1) -- (5) のいずれかが $f$ と $g$ に対して
成り立つと仮定する。このとき
$$
\deg(X/Z) = \deg(X/Y) \deg(Y/Z).
$$
\end{lemma}

\begin{proof}
これは体の有限拡大の塔における次数の乗法性から従う。Fields, Lemma
\ref{fields-lemma-multiplicativity-degrees} を参照せよ。
\end{proof}










\section{Weil 因子}
\label{spaces-over-fields-section-Weil-divisors}

\noindent
本節は Divisors, Section \ref{divisors-section-Weil-divisors} の対応物である。

\medskip\noindent
局所 Noether 整代数空間に対して Weil 因子と Weil 因子の有理同値を導入する。
ここでは代数空間が準コンパクトであると仮定しないため、Weil 因子を定義する際には
少し注意が必要である。例えば有理関数が無限個の極をもつことがあるので、
素因子の無限和を許さなければならない。準コンパクトな場合には、通常どおり
Weil 因子は有限和となる。以下の基本補題は、閉部分空間の族が局所有限であることを
示すためにしばしば用いる。

\begin{lemma}
\label{spaces-over-fields-lemma-components-locally-finite}
$S$ をスキームとし、$X$ を $S$ 上の局所 Noether 代数空間とする。
$T \subset |X|$ が閉部分集合ならば、$T$ の既約成分の族は局所有限である。
\end{lemma}

\begin{proof}
位相空間 $|X|$ は局所 Noether である（Properties of Spaces, Lemma
\ref{spaces-properties-lemma-Noetherian-topology}）。Noether 位相空間は有限個の既約成分を
もち、Noether 空間の部分空間も Noether である（Topology, Lemma
\ref{topology-lemma-Noetherian}）。従って、補題は局所有限の定義
（Topology, Definition \ref{topology-definition-locally-finite}）から従う。
\end{proof}

\noindent
$S$ をスキームとする。$X$ を $S$ 上の decent 代数空間とする。
$Z$ を $X$ の整な閉部分空間とし、$\xi \in |Z|$ を一般点とする。このとき
$|Z|$ の $|X|$ における余次元は、Decent Spaces, Lemma
\ref{decent-spaces-lemma-codimension-local-ring} により、$X$ の $\xi$ における
局所環の次元に等しい。このことを {\it $\xi$ は余次元 $1$ の点である（$X$ 上で）}
とも表すことを思い出そう。Properties of Spaces, Definition
\ref{spaces-properties-definition-dimension-local-ring} を参照せよ。

\begin{definition}
\label{spaces-over-fields-definition-Weil-divisor}
$S$ をスキームとする。$X$ を $S$ 上の局所 Noether 整代数空間とする。
\begin{enumerate}
\item {\it 素因子}とは、整な閉部分空間 $Z \subset X$ で余次元 $1$ のもの、すなわち
$|Z|$ の一般点が余次元 $1$ の $X$ の点となるものである。
\item {\it Weil 因子}とは形式和 $D = \sum n_Z Z$ であって、和は $X$ の素因子を
走り、族 $\{|Z| : n_Z \not = 0\}$ が $|X|$ において局所有限となるものをいう
（Topology, Definition \ref{topology-definition-locally-finite}）。
\end{enumerate}
$X$ 上のすべての Weil 因子からなる群を $\text{Div}(X)$ と記す。
\end{definition}

\noindent
次の課題は、有理関数に付随する Weil 因子を定義することである。そのためには、
局所 Noether 整代数空間 $X$ 上の有理関数について、素因子 $Z$ に沿う消滅次数を
定義する必要がある。$\xi \in |Z|$ を一般点とする。ここで、局所環
$\mathcal{O}_{X, \xi}$ は存在せず、Hensel 局所環 $\mathcal{O}_{X, \xi}^h$ は
整域とは限らないという問題が生じる。Example \ref{spaces-over-fields-example-not-a-domain} を参照せよ。
これを回避するため、次の補題を用いる。

\begin{lemma}
\label{spaces-over-fields-lemma-order-vanishing}
$S$ をスキームとする。$X$ を $S$ 上の局所 Noether 整代数空間とする。
$Z \subset X$ を素因子とし、$\xi \in |Z|$ を一般点とする。このとき Hensel 局所環
$\mathcal{O}_{X, \xi}^h$ は被約な $1$ 次元 Noether 局所環であり、標準的な単射
$$
R(X) \longrightarrow Q(\mathcal{O}_{X, \xi}^h)
$$
が存在する。これは関数体 $R(X)$、すなわち $X$ の関数体から全分数環への写像である。
\end{lemma}

\begin{proof}
Decent Spaces, Section
\ref{decent-spaces-section-residue-fields-henselian-local-rings} の結果を用いる。
$(U, u) \to (X, \xi)$ を基本的な \'etale 近傍とする。$U$ は局所 Noether かつ
被約であることに注意せよ。従って $\mathcal{O}_{U, u}$ は（素因子の定義により）
$1$ 次元の被約 Noether 環である。$U$ を $u$ のアフィン開近傍で置き換えることにより、
$U$ は Noether かつアフィンであると仮定してよい。さらに $U$ をより小さい開集合で
置き換え、$U$ のすべての既約成分が $u$ を通ると仮定してよい。
$U \to X$ は開写像で $X$ は既約なので、$U \to X$ は支配的である。
従って、有理写像を合成することにより環準同型 $R(X) \to R(U)$ を得る。
Morphisms of Spaces, Section \ref{spaces-morphisms-section-rational-maps} を参照せよ。
$R(X)$ は体なので、この準同型は単射である。$U$ の選び方により、$R(U)$ は全商環
$Q(\mathcal{O}_{U, u})$ である。Morphisms, Lemma
\ref{morphisms-lemma-integral-scheme-rational-functions} および Algebra, Lemma
\ref{algebra-lemma-total-ring-fractions-no-embedded-points} を参照せよ。

\medskip\noindent
ここまでで、$\mathcal{O}_{U, u}$ を
% P09-ERR-0246
$\mathcal{O}_{X, \xi}^h$ の代わりに用いた場合の補題のすべての主張を証明した。しかし
$\mathcal{O}_{X, \xi}^h$ は $\mathcal{O}_{U, u}$ の Hensel 化である。従って
$\mathcal{O}_{X, \xi}^h$ は $1$ 次元の被約 Noether 環である。More on Algebra, Lemmas
\ref{more-algebra-lemma-henselization-reduced},
\ref{more-algebra-lemma-henselization-dimension}, and
\ref{more-algebra-lemma-henselization-noetherian} を参照せよ。
More on Algebra, Lemma \ref{more-algebra-lemma-dumb-properties-henselization} により
$\mathcal{O}_{U, u} \to \mathcal{O}_{X, \xi}^h$ は忠実平坦なので、非零因子を
非零因子へ写す。従って標準的な写像
$Q(\mathcal{O}_{U, u}) \to Q(\mathcal{O}_{X, \xi}^h)$ を得て、所望の写像を得る。
この写像が
% P09-ERR-0247
$(U, u) \to (X, x)$ の選択に依存しないことの確認は省略する。少し良い方法は、まず
$\colim Q(\mathcal{O}_{U, u}) = Q(\mathcal{O}_{X, \xi}^h)$
に注意することである。
\end{proof}

\begin{definition}
\label{spaces-over-fields-definition-order-vanishing}
$S$ をスキームとする。$X$ を $S$ 上の局所 Noether 整代数空間とする。
$f \in R(X)^*$ とする。各素因子 $Z \subset X$ に対して、{\it $f$ の $Z$ に沿う
消滅次数}を整数
$$
\text{ord}_Z(f) =
\text{length}_{\mathcal{O}_{X, \xi}^h}
(\mathcal{O}_{X, \xi}^h/a \mathcal{O}_{X, \xi}^h) -
\text{length}_{\mathcal{O}_{X, \xi}^h}
(\mathcal{O}_{X, \xi}^h/b \mathcal{O}_{X, \xi}^h)
$$
として定義する。ここで $a, b \in \mathcal{O}_{X, \xi}^h$ は非零因子で、
$f$ の $Q(\mathcal{O}_{X, \xi}^h)$ における像（Lemma
\ref{spaces-over-fields-lemma-order-vanishing}）が $a/b$ に等しいものとする。
これは Algebra, Lemma \ref{algebra-lemma-ord-additive} により良定義である。
\end{definition}

\noindent
$\mathcal{O}_{X, \xi}^h$ がたまたま整域であるならば
$$
\text{ord}_Z(f) = \text{ord}_{\mathcal{O}_{X, \xi}^h}(f)
$$
を得る。ここで右辺は Algebra, Definition \ref{algebra-definition-ord} の概念である。
$f, g \in R(X)^*$ に対して
$$
\text{ord}_Z(fg) = \text{ord}_Z(f) + \text{ord}_Z(g).
$$
であることにも注意せよ。もちろん $\text{ord}_Z(f) < 0$ となることもある。
この場合、$f$ は $Z$ に沿って {\it 極} をもつといい、
$-\text{ord}_Z(f) > 0$ を {\it $f$ の $Z$ に沿う極の位数} という。
条件 $\text{ord}_Z(f) \geq 0$ は、条件 $f \in \mathcal{O}_{X, \xi}^h$ と
{\bf 同値ではない} ことが重要である。ただし、局所環
% P09-ERR-0248
$\mathcal{O}_{X, \xi}$ が離散付値環である場合を除く。

\begin{lemma}
\label{spaces-over-fields-lemma-order-vanishing-agrees}
$S$ をスキームとする。$X$ を $S$ 上の局所 Noether 整代数空間とする。
$f \in R(X)^*$ とする。素因子 $Z \subset X$ が $X$ のスキーム的軌跡と交わるならば、
Definition \ref{spaces-over-fields-definition-order-vanishing} の消滅次数 $\text{ord}_Z(f)$ は、
Divisors, Definition \ref{divisors-definition-order-vanishing} の消滅次数と一致する。
\end{lemma}

\begin{proof}
$X$ を縮小することにより、$X$ は整 Noether スキームであると仮定してよい。
$\xi \in Z$ が一般点を表すならば、$\mathcal{O}_{X, \xi}^h$ は
$\mathcal{O}_{X, \xi}$ の Hensel 化である（Decent Spaces, Lemma
\ref{decent-spaces-lemma-describe-henselian-local-ring}）。補題を証明するための必要十分条件は
$$
\text{length}_{\mathcal{O}_{X, \xi}}
(\mathcal{O}_{X, \xi}/a \mathcal{O}_{X, \xi}) =
\text{length}_{\mathcal{O}_{X, \xi}^h}
(\mathcal{O}_{X, \xi}^h/a \mathcal{O}_{X, \xi}^h)
$$
である。これは Algebra, Lemma \ref{algebra-lemma-pullback-module} から直ちに従う
（$\mathcal{O}_{X, \xi} \to \mathcal{O}_{X, \xi}^h$ が Noether 局所環の間の
平坦な局所環準同型であることも用いる）。
\end{proof}

\begin{lemma}
\label{spaces-over-fields-lemma-divisor-locally-finite}
$S$ をスキームとする。$X$ を $S$ 上の局所 Noether 整代数空間とする。
$f \in R(X)^*$ とする。このとき次の族
% P09-ERR-0249
$$
\{Z \subset X \mid Z\text{ a prime divisor with generic point }\xi
\text{ and }f\text{ not in }\mathcal{O}_{X, \xi}\}
$$
および
$$
\{Z \subset X \mid Z \text{ a prime divisor and }\text{ord}_Z(f) \not = 0\}
$$
は $X$ において局所有限である。
\end{lemma}

\begin{proof}
非空開部分空間 $U \subset X$ で、$f$ が $\Gamma(U, \mathcal{O}_X^*)$ の切断に
対応するものが存在する。従って、補題の集合に現れうる素因子はすべて
$|X| \setminus |U|$ の既約成分に対応する。ゆえに Lemma
\ref{spaces-over-fields-lemma-components-locally-finite} が所望の結果を与える。
\end{proof}

\noindent
この補題により次の定義が可能となる。

\begin{definition}
\label{spaces-over-fields-definition-principal-divisor}
$S$ をスキームとする。$X$ を $S$ 上の局所 Noether 整代数空間とする。
$f \in R(X)^*$ とする。{\it $f$ に付随する主 Weil 因子}とは Weil 因子
$$
\text{div}(f) = \text{div}_X(f) = \sum \text{ord}_Z(f) [Z]
$$
のことである。ここで和は素因子を走り、$\text{ord}_Z(f)$ は Definition
\ref{spaces-over-fields-definition-order-vanishing} のとおりである。Lemma
\ref{spaces-over-fields-lemma-divisor-locally-finite} により、これは意味をもつ。
\end{definition}

\begin{lemma}
\label{spaces-over-fields-lemma-div-additive}
$S$ をスキームとする。$X$ を $S$ 上の局所 Noether 整代数空間とする。
$f, g \in R(X)^*$ とする。このとき
$$
\text{div}_X(fg) = \text{div}_X(f) + \text{div}_X(g)
$$
が $X$ 上の Weil 因子として成り立つ。
\end{lemma}

\begin{proof}
これは $\text{ord}$ 関数の加法性から明らかである。
\end{proof}

\noindent
% P09-ERR-0250
上の補題から、主 Weil 因子の族はすべての Weil 因子からなる群の部分群をなすことが
分かる。これにより次の定義を得る。

\begin{definition}
\label{spaces-over-fields-definition-class-group}
$S$ をスキームとする。$X$ を $S$ 上の局所 Noether 整代数空間とする。
$X$ の {\it Weil 因子類群}とは、Weil 因子の群を主 Weil 因子の部分群で割った商である。
記号は $\text{Cl}(X)$ とする。
\end{definition}

\noindent
構成により正合複体
\begin{equation}
\label{spaces-over-fields-equation-Weil-divisor-class}
R(X)^* \xrightarrow{\text{div}} \text{Div}(X) \to \text{Cl}(X) \to 0
\end{equation}

これは $\text{Cl}(X)$ の表示とみなすことができる。次の課題は、Weil 因子類群を
Picard 群と関連づけることである。

\begin{example}
\label{spaces-over-fields-example-Z-mod-2}
これは Morphisms of Spaces, Example
\ref{spaces-morphisms-example-universal-homeomorphism} の続きである。代数空間
$X = \mathbf{A}^1_k/\{t \sim -t \mid t \not = 0\}$ を考える。
これは体 $k$ 上の滑らかな代数空間である。普遍的同相写像
$$
X \longrightarrow \mathbf{A}^1_k = \Spec(k[t])
$$
が存在し、$\mathbf{A}^1_k \setminus \{0\}$ 上では同型である。
従って $X$ は Noether かつ整である。$\dim(X) = 1$ なので、$X$ の素因子は
$X$ の閉点である。$x \in |X|$ を、$0 \in \mathbf{A}^1_k$ の上にある
一意な閉点とする。
% P09-ERR-0251
$X \setminus \{x\}$ は $\mathbf{A}^1 \setminus \{0\}$ へ同型に写るので、
$\text{Cl}(X)$ において $x$ と異なる閉点の類は零である。しかし $t$ の $X$ 上の
因子は $2[x]$ である。従って $\text{Cl}(X) = \mathbf{Z}/2\mathbf{Z}$ である。
\end{example}

\begin{example}
\label{spaces-over-fields-example-not-a-domain}
$k$ を体とする。
$$
U = \Spec(k[x, y]/(xy))
$$
を $\mathbf{A}^2_k$ における二本の座標軸の和とする。
% P09-ERR-0252
$\Delta : U \to U \times_k U$ を対角射とし、
$\Delta' : U \to U \times_k U$ を写像 $u \mapsto (u, \sigma(u))$ とする。
ここで $\sigma : U \to U$、$(x, y) \mapsto (y, x)$ は二本の座標軸を
入れ替える自己同型である。
$$
R = \Delta(U) \amalg \Delta'(U \setminus \{0_U\})
$$
とおく。ここで $0_U \in U$ は原点である。$R$ が $U$ 上の \'etale 同値関係である
ことは容易に分かる。商 $X = U/R$ は代数空間である。射
$U \to \mathbf{A}^1_k$、$(x, y) \mapsto x + y$ は $R$-不変なので、射
$$
X \longrightarrow \mathbf{A}^1_k
$$
を定める。この射は普遍的同相写像であり、
$\mathbf{A}^1_k \setminus \{0\}$ 上では同型である。従って $X$ は整かつ
Noether である。Example \ref{spaces-over-fields-example-Z-mod-2} とまったく同様にして、読者は
$\text{Cl}(X) = \mathbf{Z}/2\mathbf{Z}$ であり、その生成元が一意な閉点
$x \in |X|$、すなわち $0 \in \mathbf{A}^1_k$ に写るものに対応することを示せる。
しかしこの場合、$X$ の $x$ における Hensel 局所環は整域ではない。
これは $\mathcal{O}_{U, 0_U}$ の Hensel 化だからである。
\end{example}














\section{可逆加群に付随する Weil 因子類}
\label{spaces-over-fields-section-c1}

\noindent
本節では Section \ref{spaces-over-fields-section-Weil-divisors} とまったく同じ段階をたどり、
局所 Noether 整代数空間に対する標準的な写像
$\Pic(X) \to \text{Cl}(X)$ を定義する。

\medskip\noindent
$S$ をスキームとする。$X$ を $S$ 上の局所 Noether 整代数空間とする。
$\mathcal{L}$ を可逆 $\mathcal{O}_X$-加群とする。Divisors on Spaces, Lemma
\ref{spaces-divisors-lemma-regular-meromorphic-section-exists} により、
正則有理型切断 $s \in \Gamma(X, \mathcal{K}_X(\mathcal{L}))$ が存在する。
実際、Divisors on Spaces, Lemma
\ref{spaces-divisors-lemma-compute-meromorphic} により、これは
$\mathcal{L}_\eta$ の非零元と同じものである。ここで $\eta \in |X|$ は一般点である。
同じ補題によれば、$\mathcal{L} = \mathcal{O}_X$ ならば $s$ は $X$ 上の非零有理関数と
同じものである（従って以下の構成は Section
\ref{spaces-over-fields-section-Weil-divisors} の構成と一致する）。

\medskip\noindent
$Z \subset X$ を素因子とし、$\xi \in |Z|$ を一般点とする。
$s$ の $Z$ に沿う消滅次数を定義する。標準的な射
$$
c_\xi : \Spec(\mathcal{O}_{X, \xi}^h) \longrightarrow X
$$
を考える。その始域は $X$ の $\xi$ における Hensel 局所環のスペクトルである
（Decent Spaces, Definition
% P09-ERR-0253
\ref{decent-spaces-definition-henselian-local-ring}）。
引き戻し $\mathcal{L}_\xi = c_\xi^*\mathcal{L}$ は可逆加群なので自明である。
生成元 $s_\xi$ を $\mathcal{L}_\xi$ に対して選ぶ。$c_\xi$ は平坦なので、有理型関数および
（正則）切断の引き戻しは $c_\xi$ に対して定義される。Divisors on Spaces,
Definition \ref{spaces-divisors-definition-pullback-meromorphic-sections} および
Lemmas \ref{spaces-divisors-lemma-pullback-meromorphic-sections-defined} and
\ref{spaces-divisors-lemma-meromorphic-sections-pullback} を参照せよ。従って
$$
c_\xi^*(s) = f s_\xi
$$
を得る。ここで $f \in Q(\mathcal{O}_{X, \xi}^h)$ は非零因子である。
ここでは Divisors, Lemma \ref{divisors-lemma-locally-Noetherian-K} を用い、
$\mathcal{L}_\xi \cong \mathcal{O}_{\Spec(\mathcal{O}_{X, \xi}^h)}$ の
有理型切断の空間を $\mathcal{O}_{X, \xi}^h$ の全分数環によって同一視している。
この元を
$$
s/s_\xi = f \in  Q(\mathcal{O}_{X, \xi}^h)
$$
と書くことにする。$f = s/s_\xi$ は $uf$ に置き換わることに注意せよ。
ここで $u \in \mathcal{O}_{X, \xi}^h$ は単元であり、これは $s_\xi$ の選び方を
変えた場合に起こる。

\begin{definition}
\label{spaces-over-fields-definition-order-vanishing-meromorphic}
$S$ をスキームとする。
% P09-ERR-0254
$X$ を $S$ 上の局所 Noether 整代数空間とする。$\mathcal{L}$ を可逆
$\mathcal{O}_X$-加群とする。$s \in \Gamma(X, \mathcal{K}_X(\mathcal{L}))$ を
$\mathcal{L}$ の正則有理型切断とする。各素因子 $Z \subset X$ で一般点
$\xi \in |Z|$ をもつものに対して、{\it $s$ の $Z$ に沿う消滅次数}を整数
$$
\text{ord}_{Z, \mathcal{L}}(s) =
\text{length}_{\mathcal{O}_{X, \xi}^h}
(\mathcal{O}_{X, \xi}^h/a \mathcal{O}_{X, \xi}^h) -
\text{length}_{\mathcal{O}_{X, \xi}^h}
(\mathcal{O}_{X, \xi}^h/b \mathcal{O}_{X, \xi}^h)
$$
として定義する。ここで $a, b \in \mathcal{O}_{X, \xi}^h$ は非零因子であり、
上で構成した $s/s_\xi$ という $Q(\mathcal{O}_{X, \xi}^h)$ の元が
$a/b$ に等しいものとする。上の議論と Algebra, Lemma
\ref{algebra-lemma-ord-additive} により、これは良定義である。
\end{definition}

\noindent
上で説明したように、正則有理型切断 $s$ が $\mathcal{O}_X$ のものであるならば、
$s = f \cdot 1$ と書ける。ここで $f$ は $X$ 上の非零有理関数であり、
$\text{ord}_Z(f) = \text{ord}_{Z, \mathcal{O}_X}(s)$ が成り立つ。
主因子の場合と同様、次の補題が成り立つ。

\begin{lemma}
\label{spaces-over-fields-lemma-divisor-meromorphic-locally-finite}
$S$ をスキームとする。$X$ を $S$ 上の局所 Noether 整代数空間とする。
$\mathcal{L}$ を可逆 $\mathcal{O}_X$-加群とする。
% P09-ERR-0255
$s \in \mathcal{K}_X(\mathcal{L})$ を $\mathcal{L}$ の正則（すなわち非零）有理型切断とする。
このとき集合
$$
\{Z \subset X \mid Z\text{ は一般点 }\xi
\text{ をもつ素因子で、かつ }s\text{ が属さない }\mathcal{L}_\xi\}
$$
および
$$
\{Z \subset X \mid Z \text{ は素因子であり }
\text{ord}_{Z, \mathcal{L}}(s) \not = 0\}
$$
は $X$ において局所有限である。
\end{lemma}

\begin{proof}
非空開部分空間 $U \subset X$ で、$s$ が $\Gamma(U, \mathcal{L})$ の切断に対応し、
その切断が $\mathcal{L}$ を $U$ 上で生成するものが存在する。従って、補題の集合に
現れうる素因子はすべて $|X| \setminus |U|$ の既約成分に対応する。従って Lemma
% P09-ERR-0256
\ref{spaces-over-fields-lemma-components-locally-finite} が所望の結果を与える。
\end{proof}

\begin{lemma}
\label{spaces-over-fields-lemma-divisor-meromorphic-well-defined}
$S$ をスキームとする。
% P09-ERR-0257
$X$ を $S$ 上の局所 Noether 整代数空間とする。$\mathcal{L}$ を可逆
$\mathcal{O}_X$-加群とする。$s, s' \in \mathcal{K}_X(\mathcal{L})$ を
$\mathcal{L}$ の非零有理型切断とする。このとき $f = s/s'$ は
$R(X)^*$ の元であり、
$$
\sum \text{ord}_{Z, \mathcal{L}}(s)[Z]
=
\sum \text{ord}_{Z, \mathcal{L}}(s')[Z]
+
\text{div}(f)
$$
が Weil 因子として成り立つ。
\end{lemma}

\begin{proof}
これは定義から明らかである。Lemma
\ref{spaces-over-fields-lemma-divisor-meromorphic-locally-finite} によって、これらの和が実際に
Weil 因子となることが保証される点に注意せよ。
\end{proof}

\begin{definition}
\label{spaces-over-fields-definition-divisor-invertible-sheaf}
$S$ をスキームとする。$X$ を $S$ 上の局所 Noether 整代数空間とする。
$\mathcal{L}$ を可逆 $\mathcal{O}_X$-加群とする。
\begin{enumerate}
\item 任意の非零有理型切断 $s$ で $\mathcal{L}$ のものに対し、{\it $s$ に付随する
Weil 因子}を
$$
\text{div}_\mathcal{L}(s) =
\sum \text{ord}_{Z, \mathcal{L}}(s) [Z] \in \text{Div}(X)
$$
と定義する。ここで和は素因子を走る。Lemma
\ref{spaces-over-fields-lemma-divisor-meromorphic-locally-finite} により、これは良定義である。
\item {\it $\mathcal{L}$ に付随する Weil 因子類}を
$\text{div}_\mathcal{L}(s)$ の $\text{Cl}(X)$ における像として定義する。
ここで $s$ は $\mathcal{L}$ の $X$ 上の任意の非零有理型切断である。Lemma
\ref{spaces-over-fields-lemma-divisor-meromorphic-well-defined} により、これは良定義である。
\end{enumerate}
\end{definition}

\noindent
予想どおり、この構成は可逆加群について加法的である。

\begin{lemma}
\label{spaces-over-fields-lemma-c1-additive}
$S$ をスキームとする。$X$ を $S$ 上の局所 Noether 整代数空間とする。
$\mathcal{L}$、$\mathcal{N}$ を可逆 $\mathcal{O}_X$-加群とする。
$s$、それぞれ $t$ を $\mathcal{L}$、それぞれ $\mathcal{N}$ の非零有理型切断とする。
このとき $st$ は $\mathcal{L} \otimes_{\mathcal{O}_X} \mathcal{N}$ の
非零有理型切断であり、
$$
\text{div}_{\mathcal{L} \otimes \mathcal{N}}(st)
=
\text{div}_\mathcal{L}(s) + \text{div}_\mathcal{N}(t)
$$
が $\text{Div}(X)$ において成り立つ。特に
$\mathcal{L} \otimes_{\mathcal{O}_X} \mathcal{N}$ の Weil 因子類は、
$\mathcal{L}$ と $\mathcal{N}$ の Weil 因子類の和である。
\end{lemma}

\begin{proof}
$s$、それぞれ $t$ を $\mathcal{L}$、それぞれ $\mathcal{N}$ の非零有理型切断とする。
このとき $st$ は $\mathcal{L} \otimes \mathcal{N}$ の非零有理型切断である。
$Z \subset X$ を素因子とし、$\xi \in |Z|$ をその一般点とする。本節で先に導入した
記法を用いて、生成元 $s_\xi \in \mathcal{L}_\xi$ と
$t_\xi \in \mathcal{N}_\xi$ を選ぶ。このとき $s_\xi \otimes t_\xi$ は
$(\mathcal{L} \otimes \mathcal{N})_\xi$ の生成元である。
$st/(s_\xi t_\xi) = (s/s_\xi)(t/t_\xi)$ が
$Q(\mathcal{O}_{X, \xi}^h)$ において成り立つ。Algebra, Lemma
\ref{algebra-lemma-ord-additive} の加法性を適用すると
$$
\text{div}_{\mathcal{L} \otimes \mathcal{N}, Z}(st)
=
\text{div}_{\mathcal{L}, Z}(s) + \text{div}_{\mathcal{N}, Z}(t)
$$
を得る。
% P09-ERR-0258
いくつかの詳細は省略する。
\end{proof}

\noindent
$S$ をスキームとする。$X$ を $S$ 上の局所 Noether 整代数空間とする。
上の構成と補題により、Abel 群の準同型
\begin{equation}
\label{spaces-over-fields-equation-c1}
\Pic(X) \longrightarrow \text{Cl}(X)
\end{equation}
を得る。これは可逆加群をその Weil 因子類へ写す。

\begin{lemma}
\label{spaces-over-fields-lemma-normal-c1-injective}
$S$ をスキームとする。$X$ を $S$ 上の局所 Noether 整代数空間とする。
$X$ が正規ならば、写像 (\ref{spaces-over-fields-equation-c1})
$\Pic(X) \to \text{Cl}(X)$ は単射である。
\end{lemma}

\begin{proof}
$\mathcal{L}$ を可逆 $\mathcal{O}_X$-加群とし、それに付随する Weil 因子類が
自明であるとする。$s$ を $\mathcal{L}$ の正則有理型切断とする。この仮定は、
$\text{div}_\mathcal{L}(s) = \text{div}(f)$ が、ある
$f \in R(X)^*$ に対して成り立つことを意味する。
すると $t = f^{-1}s$ は $\mathcal{L}$ の正則有理型切断で、
$\text{div}_\mathcal{L}(t) = 0$ を満たす。Lemma
\ref{spaces-over-fields-lemma-divisor-meromorphic-well-defined} を参照せよ。
$t$ が $\mathcal{L}$ の自明化を定めると主張する。この主張で補題の証明は終わる。
まだ多くの理論を自由に使えないため、この主張の証明は少し不格好である。読者には
% P09-ERR-0259
証明を飛ばすことを勧める。

\medskip\noindent
この主張は \'etale 局所的に確認できる。$U \in X_\etale$ をアフィンとし、
$\mathcal{L}|_U$ が自明であるとする。$s_U \in \Gamma(U, \mathcal{L}|_U)$ を
一つの自明化とする。Properties, Lemma
\ref{properties-lemma-normal-locally-finite-nr-irreducibles} により、
$U$ も整であると仮定してよい。$U = \Spec(A)$ を正規 Noether 整域 $A$ の
スペクトルとして書き、その分数体を $K$ とする。
$t|_U = f s_U$ と書ける。ここで $f$ は $K$ の元である。例えば Divisors on Spaces, Lemma
\ref{spaces-divisors-lemma-meromorphic-quasi-coherent} を参照せよ。
$\mathfrak p \subset A$ を高さ一の素イデアルとし、余次元 $1$ の点
$u \in U$ に対応するものとする。この点は余次元 $1$ の点
$\xi \in |X|$ へ写る。本節の初めと同様に自明化
$s_\xi$ を $\mathcal{L}_\xi$ に対して選ぶ。幾何学的点 $\overline{u}$ を
$U$ から選び、$u$ の上にあるものとする。
このとき
$$
(\mathcal{O}_{X, \xi}^h)^{sh} =
\mathcal{O}_{X, \overline{u}} =
\mathcal{O}_{U, u}^{sh} = (A_\mathfrak p)^{sh}
$$
である。Decent Spaces, Lemmas
\ref{decent-spaces-lemma-henselian-local-ring-strict} および
Properties of Spaces, Lemma
\ref{spaces-properties-lemma-describe-etale-local-ring} を参照せよ。
$X$ の正規性から、これらはすべて離散付値環である。
自明化 $s_U$ と $s_\xi$ は、$\mathcal{L}$ を
$\Spec(\mathcal{O}_{X, \overline{u}})$ へ引き戻した切断として、単元だけ異なる。
$t = f_\xi s_\xi$ と書く。ここで
$f_\xi \in Q(\mathcal{O}_{X, \xi}^h)$ である。従って $f_\xi$ と $f$ は
$Q(\mathcal{O}_{X, \overline{u}})$ における単元だけ異なる。
$Z \subset X$ が $\xi$ に対応する素因子を表すならば（Lemma
\ref{spaces-over-fields-lemma-decent-irreducible-closed}）、
$0 = \text{ord}_{Z, \mathcal{L}}(t) =
\text{ord}_{\mathcal{O}_{X, \xi}^h}(f_\xi)$
である。また $\mathcal{O}_{X, \xi}^h$ は離散付値環なので $f_\xi$ は単元である。
従って $f$ は $\mathcal{O}_{X, \overline{u}}$ の単元であり、特に
$f \in A_\mathfrak p$ である。Algebra, Lemma
\ref{algebra-lemma-normal-domain-intersection-localizations-height-1} により、
これは $f \in A$ を含意する。従って $t \in \Gamma(X, \mathcal{L})$ である。
$t^{-1}$ を $\mathcal{L}^{\otimes -1}$ の有理型切断とみなし、同じ議論を繰り返せば
証明が完了する。
\end{proof}


















\section{修正とオルタレーション}
\label{spaces-over-fields-section-modifications-alterations}

\noindent
整代数空間の概念を用いて、修正を次のように定義できる。

\begin{definition}
\label{spaces-over-fields-definition-modification}
$S$ をスキームとする。$X$ を $S$ 上の整代数空間とする。{\it $X$ の修正}とは、
双有理固有射 $f : X' \to X$ で、$S$ 上の代数空間の射であり、$X'$ が整であるものをいう。
\end{definition}

\noindent
代数空間の双有理射については Decent Spaces, Definition
\ref{decent-spaces-definition-birational} を参照せよ。

\begin{lemma}
\label{spaces-over-fields-lemma-modification-iso-over-open}
$f : X' \to X$ を Definition \ref{spaces-over-fields-definition-modification} の修正とする。
非空開集合 $U \subset X$ で、$f^{-1}(U) \to U$ が同型となるものが存在する。
\end{lemma}

\begin{proof}
Lemma \ref{spaces-over-fields-lemma-finite-degree} により、非空 $U \subset X$ で
$f^{-1}(U) \to U$ が有限となるものが存在する。一般平坦性
（Morphisms of Spaces, Proposition
\ref{spaces-morphisms-proposition-generic-flatness-reduced}）により、
$f^{-1}(U) \to U$ は平坦かつ有限表示であると仮定してよい。従って
$f^{-1}(U) \to U$ は有限局所自由である（Morphisms of Spaces, Lemma
\ref{spaces-morphisms-lemma-finite-flat}）。$f$ は双有理なので、$X'$ の $X$ 上の
次数は $1$ である。従って $f^{-1}(U) \to U$ は次数 $1$ の有限局所自由射、
すなわち同型である。
\end{proof}

\begin{definition}
\label{spaces-over-fields-definition-alteration}
$S$ をスキームとする。$X$ を $S$ 上の整代数空間とする。{\it $X$ の
オルタレーション}とは、固有支配射 $f : Y \to X$ で、$S$ 上の代数空間の射であり、
$Y$ が整で、$f^{-1}(U) \to U$ が有限となる非空開集合 $U \subset X$ が
存在するものをいう。
\end{definition}

\noindent
$f : Y \to X$ が整代数空間の間の支配固有射ならば、一般点において誘導される
剰余体の拡大が有限であるとき、これはオルタレーションである。正確には次が成り立つ。
% P09-ERR-0260

\begin{lemma}
\label{spaces-over-fields-lemma-alteration-generically-finite}
$S$ をスキームとする。$f : X \to Y$ を $S$ 上の整代数空間の固有支配射とする。
$f$ がオルタレーションであることと、Lemma \ref{spaces-over-fields-lemma-finite-degree} の同値な条件
(1) -- (6) のいずれかが成り立つことは同値である。
\end{lemma}

\begin{proof}
これは本文で参照した補題の直接の帰結である。
\end{proof}

\begin{lemma}
\label{spaces-over-fields-lemma-alteration-contained-in}
$S$ をスキームとする。$f : X \to Y$ を $S$ 上の代数空間の固有全射とする。
$Y$ が整であると仮定する。このとき整な閉部分空間 $X' \subset X$ で、
$f' = f|_{X'} : X' \to Y$ がオルタレーションとなるものが存在する。
\end{lemma}

\begin{proof}
$V \subset Y$ を非空アフィン開集合とする（Decent Spaces, Theorem
\ref{decent-spaces-theorem-decent-open-dense-scheme}）。
$\eta \in V$ を一般点とする。このとき $X_\eta$ は $\eta$ 上の非空固有代数空間である。
閉点 $x \in |X_\eta|$ を選ぶ（これは $|X_\eta|$ が準コンパクトな sober 位相空間
だから存在する。Decent Spaces, Proposition
\ref{decent-spaces-proposition-reasonable-sober} および Topology, Lemma
\ref{topology-lemma-quasi-compact-closed-point} を参照）。
% P09-ERR-0261
$X'$ を $\overline{\{x\}} \subset |X|$ 上の被約誘導閉部分空間構造とする
（Properties of Spaces, Definition
\ref{spaces-properties-definition-reduced-induced-space}）。
% P09-ERR-0262
$f' : X' \to Y$ は、その像が $\eta$ を含むため全射である。また $f'$ は閉埋め込みと
固有射の合成なので固有である。最後に、ファイバー $X'_\eta$ は一点だけをもつ。
これを見るには $X \to Y$ と $X' \to Y$ の双方、および点 $\eta$ に
Decent Spaces, Lemma \ref{decent-spaces-lemma-topology-fibre} を用いればよい。
$Y$ は decent で $X' \to Y$ は分離的なので、$X'$ は decent である
（Decent Spaces, Lemmas
% P09-ERR-0263
\ref{decent-spaces-lemma-properties-trivial-implications} and
\ref{decent-spaces-lemma-property-over-property}）。
従って Lemma \ref{spaces-over-fields-lemma-alteration-generically-finite} により、
$f'$ はオルタレーションである。
\end{proof}









\section{スキーム的軌跡}
\label{spaces-over-fields-section-schematic}

\noindent
代数空間のスキーム的軌跡については、すでにいくつもの結果を証明した。
参照箇所を列挙する。
\begin{enumerate}
\item Properties of Spaces, Sections
\ref{spaces-properties-section-schematic} and
\ref{spaces-properties-section-getting-a-scheme},
\item Decent Spaces, Section \ref{decent-spaces-section-schematic},
\item Properties of Spaces, Lemma
\ref{spaces-properties-lemma-point-like-spaces}
$\Leftarrow$
Decent Spaces, Lemma \ref{decent-spaces-lemma-decent-point-like-spaces}
$\Leftarrow$
Decent Spaces, Lemma \ref{decent-spaces-lemma-when-field},
\item Limits of Spaces, Section \ref{spaces-limits-section-affine}, and
\item Limits of Spaces, Section \ref{spaces-limits-section-representable}.
\end{enumerate}
代数空間のある種の射が自動的に表現可能となる場合がある。例えば、分離的な
局所準有限射（Morphisms of Spaces, Lemma
\ref{spaces-morphisms-lemma-locally-quasi-finite-separated-representable}）や、
平坦単射（More on Morphisms of Spaces, Lemma
\ref{spaces-more-morphisms-lemma-flat-case}）である。Section
\ref{spaces-over-fields-section-schematic-and-field-extension} では、基礎体の拡大のもとで
スキーム的軌跡がどう変化するかを調べる。

\begin{lemma}
\label{spaces-over-fields-lemma-locally-finite-type-dim-zero}
$S$ をスキームとする。$X$ を $S$ 上の代数空間とする。$X$ が次の条件の少なくとも
一つを満たすと仮定する。
\begin{enumerate}
\item $X$ は準分離で $\dim(X) = 0$ である。
\item $X$ は体 $k$ 上局所有限型で $\dim(X) = 0$ である。
\item $X$ は Noether で $\dim(X) = 0$ である。
\item ここにさらに追加する。
% P09-ERR-0264
\end{enumerate}
このとき $X$ は分離スキームであり、$X$ の任意の準コンパクト開集合はアフィンである。
\end{lemma}

\begin{proof}
$X$ の任意の準コンパクト開集合がアフィンであることを証明すれば、$X$ は分離スキーム
となる。従って $X$ は準コンパクトであると仮定してよく、$X$ がアフィンであることを
示す。場合 (2) と (3) は場合 (1) から直ちに従うが、それぞれの証明は使う理論が
著しく少ないので、(2) と (3) を別々に証明する。
% P09-ERR-0265

\medskip\noindent
場合 (3) の証明。$U$ をアフィンスキームとし、$U \to X$ を \'etale 射とする。
$R = U \times_X U$ とおく。二つの射影 $s, t : R \to U$ はスキームの
\'etale 射である。Properties of Spaces, Definition
\ref{spaces-properties-definition-dimension} により、
$\dim(U) = 0$ かつ $\dim(R) = 0$ である。$R$ は次元 $0$ の局所 Noether スキーム
なので、$R$ は Artin 局所環のスペクトルの非交和である（Properties, Lemma
\ref{properties-lemma-locally-Noetherian-dimension-0}）。
$X$ は Noether（従って準分離）と仮定したので、$R$ は準コンパクトである。
従って $R$ はアフィンスキームである（Schemes, Lemma
\ref{schemes-lemma-disjoint-union-affines} を用いる）。
\'etale 射 $s, t : R \to U$ は有限な剰余体拡大を誘導する。従って Algebra, Lemma
\ref{algebra-lemma-essentially-of-finite-type-into-artinian-local}
により（小さな詳細を一つ省略する）、$s$ と $t$ は有限である。ゆえに
Groupoids, Proposition \ref{groupoids-proposition-finite-flat-equivalence} から、
$X = U/R$ はアフィンスキームである。

\medskip\noindent
場合 (2) の証明は場合 (3) とほぼ同じである。$U$ をアフィンスキームとし、
$U \to X$ を全射な \'etale 射とする。$R = U \times_X U$ とおく。
二つの射影 $s, t : R \to U$ はスキームの \'etale 射である。
Properties of Spaces, Definition \ref{spaces-properties-definition-dimension} により、
$\dim(U) = 0$ であり、同様に $\dim(R) = 0$ である。一方、射
$U \to \Spec(k)$ は \'etale 射 $U \to X$ と $X \to \Spec(k)$ の合成なので
局所有限型である。Morphisms of Spaces, Lemmas
\ref{spaces-morphisms-lemma-composition-finite-type} and
\ref{spaces-morphisms-lemma-etale-locally-finite-type} を参照せよ。
同様に $R \to \Spec(k)$ は局所有限型である。従って Varieties, Lemma
\ref{varieties-lemma-algebraic-scheme-dim-0} により、$U$ と $R$ は
$k$ 上有限な局所 Artin $k$-代数のスペクトルの非交和である。従って
$U \times_{\Spec(k)} U$ についても同じことが成り立つ。
$$
R = U \times_X U \longrightarrow U \times_{\Spec(k)} U
$$
は単射なので、$R$ は有限個（！）の有限 $k$-代数のスペクトルの和である。
従って $R$ はアフィンである。Schemes, Lemma
\ref{schemes-lemma-disjoint-union-affines} を参照せよ。
Varieties, Lemma \ref{varieties-lemma-algebraic-scheme-dim-0} をもう一度用いると、
$R$ は $k$ 上有限である。従って $s, t$ は有限である。Morphisms, Lemma
\ref{morphisms-lemma-finite-permanence} を参照せよ。ゆえに
Groupoids, Proposition \ref{groupoids-proposition-finite-flat-equivalence} から、
$X = U/R$ はアフィンスキームである。

\medskip\noindent
場合 (1) のコホモロジー的証明。Cohomology of Spaces, Lemma
\ref{spaces-cohomology-lemma-vanishing-above-dimension} により、すべての
準連接層 $\mathcal{F}$（$X$ 上）について高次コホモロジー群が消滅する。従って
Cohomology of Spaces, Proposition
\ref{spaces-cohomology-proposition-vanishing-affine} により、$X$ はアフィン
（特にスキーム）である。

\medskip\noindent
場合 (1) の幾何学的証明。層別化
$$
\emptyset = U_{n + 1} \subset
U_n \subset U_{n - 1} \subset \ldots \subset U_1 = X
$$
と、Decent Spaces, Lemma
\ref{decent-spaces-lemma-filter-quasi-compact-quasi-separated} のような
\'etale 射 $f_p : V_p \to U_p$ を選ぶ（以下ではそのすべての性質を用いる）。
代数空間の次元の定義により $\dim(V_p) = 0$ である。従って Properties, Lemma
\ref{properties-lemma-dimension-zero} を各 $V_p$ に適用できる。
$f_p^{-1}(U_{p + 1}) \subset V_p$ は準コンパクト開集合なので、アフィンであり、
かつ閉でもある。従って $|T_p| \subset |U_p|$（前掲箇所を参照）は開かつ閉である。
ゆえに $X$ は開閉部分空間の非交和であり、それらの被約構造はスキームである。
従って $X$ はスキームである（Limits of Spaces, Lemma
\ref{spaces-limits-lemma-reduction-scheme}）。最後に、上で参照したスキームの場合を
用いれば証明が終わる。
\end{proof}

\noindent
次の補題は、準分離代数空間が余次元 $1$ を除いてスキームであることを述べる。

\begin{lemma}
\label{spaces-over-fields-lemma-generic-point-in-schematic-locus}
$S$ をスキームとする。$X$ を $S$ 上の準分離代数空間とする。$x \in |X|$ とする。
次の条件は同値である。
\begin{enumerate}
\item $x$ は余次元 $0$ の点である（$X$ 上）。
\item $X$ の $x$ における局所環の次元は $0$ である。
\item $x$ は $|X|$ の既約成分の一般点である。
\end{enumerate}
これらが成り立つならば、$X$ の開部分空間で $x$ を含み、スキームとなるものが存在する。
\end{lemma}

\begin{proof}
(1)、(2)、(3) の同値性は Decent Spaces, Lemma
\ref{decent-spaces-lemma-decent-generic-points} と、準分離代数空間は decent である
という事実（Decent Spaces, Section
\ref{decent-spaces-section-reasonable-decent}）から従う。ただし、次の段落で
この同値性のより初等的な証明を与える。

\medskip\noindent
(1) と (2) は定義により同値である（Properties of Spaces, Definition
\ref{spaces-properties-definition-dimension-local-ring}）。
(1) と (3) の同値性を証明するため、$X$ は準コンパクトであると仮定してよい。
$$
\emptyset = U_{n + 1} \subset
U_n \subset U_{n - 1} \subset \ldots \subset U_1 = X
$$
と、Decent Spaces, Lemma
\ref{decent-spaces-lemma-filter-quasi-compact-quasi-separated} のような
$f_i : V_i \to U_i$ を選ぶ。$x \in U_i$、$x \not \in U_{i + 1}$ とする。
このとき $x = f_i(y)$ が一意な $y \in V_i$ に対して成り立つ。(1) が成り立つならば、
$y$ は $V_i$ の既約成分の一般点である（Properties of Spaces, Lemma
\ref{spaces-properties-lemma-codimension-0-points}）。
$f_i^{-1}(U_{i + 1})$ は $V_i$ の準コンパクト開集合で $y$ を含まないので、
$W \subset V_i$ という $y$ の開近傍で $f_i^{-1}(V_i)$ と交わらないものが存在する
（Properties, Lemma
% P09-ERR-0266
\ref{properties-lemma-generic-point-in-constructible}、またはより簡単に Algebra, Lemma
\ref{algebra-lemma-standard-open-containing-maximal-point} を参照）。
すると $f_i|_W : W \to X$ はその像への同型であり、従って
$x = f_i(y)$ は $|X|$ の一般点である。逆に (3) が成り立つと仮定する。
$f_i$ は $\overline{\{y\}}$ を既約成分
$\overline{\{x\}}$、すなわち $|U_i|$ のものの上へ写す。
$|f_i|$ は $\overline{\{x\}}$ 上で全単射なので、
$\overline{\{y\}}$ は $U_i$ の既約成分である。従って $x$ は余次元 $0$ の点である。
% P09-ERR-0267

\medskip\noindent
補題の最後の主張は Properties of Spaces, Proposition
\ref{spaces-properties-proposition-locally-quasi-separated-open-dense-scheme} である。
\end{proof}

\noindent
次の補題は、分離的な局所 Noether 代数空間が余次元 $1$ において、すなわち
余次元 $2$ を除いてスキームであることを述べる。

\begin{lemma}
\label{spaces-over-fields-lemma-codim-1-point-in-schematic-locus}
\begin{slogan}
分離代数空間は余次元 1 においてスキームである。
\end{slogan}
$S$ をスキームとする。$X$ を $S$ 上の代数空間とする。$x \in |X|$ とする。
$X$ が分離的かつ局所 Noether であり、$X$ の $x$ における局所環の次元が
$\leq 1$ ならば（Properties of Spaces, Definition
\ref{spaces-properties-definition-dimension-local-ring}）、$X$ の開部分空間で
$x$ を含み、スキームとなるものが存在する。
\end{lemma}

\begin{proof}
（有限亜群に関する内容を避ける別の方法については、下の注意を参照せよ。）
$X$ を $x$ の準コンパクト近傍で置き換えられるので、$X$ は準コンパクト、分離的、
かつ Noether であると仮定してよい。
% P09-ERR-0268
スキーム $U$ と有限全射 $U \to X$ が存在する。Limits of Spaces, Proposition
\ref{spaces-limits-proposition-there-is-a-scheme-finite-over} を参照せよ。
$R = U \times_X U$ とおく。このとき $j : R \to U \times_S U$ は同値関係であり、
亜群スキーム $(U, R, s, t, c)$ を得る。これは $S$ 上であり、$s, t$ は有限で、
$U$ は Noether かつ分離的である。点の集合
$\{u_1, \ldots, u_n\} \subset U$ を、$x$ に写るもの全体とする。このとき Decent Spaces, Lemma
\ref{decent-spaces-lemma-dimension-local-ring-quasi-finite} により
$\dim(\mathcal{O}_{U, u_i}) \leq 1$ である。

\medskip\noindent
More on Groupoids, Lemma
\ref{more-groupoids-lemma-find-affine-codimension-1} により、$R$-不変な
アフィン開集合 $W \subset U$ で、軌道 $\{u_1, \ldots, u_n\}$ を含むものが存在する。
$U \to X$ は有限全射なので、連続写像 $|U| \to |X|$ は閉全射である。従って
Topology, Lemma \ref{topology-lemma-closed-morphism-quotient-topology} により
商写像である。ゆえに $f(W)$ は開であり、開部分空間 $X' \subset X$ で
$f : W \to X'$ が有限全射となるものが存在する。
% P09-ERR-0269
Cohomology of Spaces, Lemma
\ref{spaces-cohomology-lemma-image-affine-finite-morphism-affine-Noetherian} により
$X'$ はアフィンスキームであり、証明が終わる。
\end{proof}

\begin{remark}
\label{spaces-over-fields-remark-alternate-proof-scheme-codim-1}
Lemma \ref{spaces-over-fields-lemma-codim-1-point-in-schematic-locus} の証明の概略を与える。この証明は
More on Groupoids, Lemma
\ref{more-groupoids-lemma-find-affine-codimension-1} を用いない。

\medskip\noindent
ステップ 1。$X$ は被約 Noether 分離代数空間であると仮定してよい
（例えば Cohomology of Spaces, Lemma
\ref{spaces-cohomology-lemma-image-affine-finite-morphism-affine-Noetherian} または
Limits of Spaces, Lemma \ref{spaces-limits-lemma-reduction-scheme} による）。
また、有限全射 $Y \to X$ で、$Y$ が Noether スキームとなるものを選べる
（Limits of Spaces, Proposition
\ref{spaces-limits-proposition-there-is-a-scheme-finite-over} による）。

\medskip\noindent
ステップ 2。$X$ を $x$ の開近傍で置き換えると、双有理有限射 $X' \to X$ と
閉部分スキーム $Y' \subset X' \times_X Y$ で、$Y' \to X'$ が有限局所自由全射となる
ものが存在する。実際、$X$ は被約なので、稠密開部分空間 $U \subset X$ で
$Y$ がその上で平坦となるものが存在する（Morphisms of Spaces, Proposition
\ref{spaces-morphisms-proposition-generic-flatness-reduced}）。
このとき $U$-許容爆発 $b : \tilde X \to X$ を、狭義変換
$\tilde Y$（$Y$ のもの）が $\tilde X$ 上で平坦となるように選べる。More on Morphisms of Spaces,
Lemma \ref{spaces-more-morphisms-lemma-flat-after-blowing-up} を参照せよ。
（爆発に関する結果を避けたければ、代わりに Hilbert スキームを用いることもできる。）
次に $X' \subset \tilde X$ を $b^{-1}(U)$ のスキーム論的閉包とし、
$Y' = X' \times_{\tilde X} \tilde Y$ とおく。$x$ は余次元 $1$ の点なので、
$X' \to X$ は $x$ の近傍上で有限である（Lemma
\ref{spaces-over-fields-lemma-finite-in-codim-1}）。

\medskip\noindent
ステップ 3。$X$ を $x$ のさらに小さい近傍へ縮小すると、$X'$ はスキームとなる。
これは $Y'$ がスキームで、$Y' \to X'$ が有限局所自由であり、さらに余次元 $1$ の
点の任意の有限集合（$Y'$ のもの）がアフィン開集合に含まれるからである。
Properties of Spaces, Proposition
% P09-ERR-0270
\ref{spaces-properties-proposition-finite-flat-equivalence-global} および
Varieties, Proposition
\ref{varieties-proposition-finite-set-of-points-of-codim-1-in-affine} を用いよ。

\medskip\noindent
ステップ 4。アフィン開集合 $W' \subset X'$ で、$x$ 上にあるすべての点を含み、
$X$ の開部分空間の逆像となるものが存在する。これを証明するため、
$Z \subset X$ を $X' \to X$ が同型でない点の集合の閉包とする。
$x \in Z$ と仮定してよい。そうでなければすでに証明は終わっている。
このとき $x$ は $Z$ の既約成分の一般点であり、$X$ を縮小することにより
$Z$ はアフィンスキームであると仮定してよい（Lemma
\ref{spaces-over-fields-lemma-generic-point-in-schematic-locus}）。すると逆像 $Z' \subset X'$ も
アフィンスキームである。点 $x_1, \ldots, x_n \in Z'$ を、$x$ に写るもの全体とする。
アフィン開集合 $W'$ を $X'$ の中に見つけられ、その $Z'$ との交わりは
$Z$ の主開集合で $x$ を含むものの逆像となる。実際、まず Varieties, Proposition
\ref{varieties-proposition-finite-set-of-points-of-codim-1-in-affine} により、
アフィン開集合 $W' \subset X'$ を選び、$x_1, \ldots, x_n$ を含むものとする。
次に主開集合 $D(f) \subset Z$ で $x$ を含むものを、その逆像
$D(f|_{Z'})$ が $W' \cap Z'$ に含まれるように選ぶ。
さらに $f' \in \Gamma(W', \mathcal{O}_{W'})$ を選び、その制限を $f|_{Z'}$ とし、
$W'$ を $D(f') \subset W'$ で置き換える。$X' \to X$ は
$Z' \to Z$ の外で同型なので、この $W'$ の選び方により、その像
$W \subset X$（$W'$ の像）は開で、その逆像 $W'$ は $X'$ にある。

\medskip\noindent
ステップ 5。このとき $W' \to W$ は有限全射であり、Cohomology of Spaces, Lemma
\ref{spaces-cohomology-lemma-image-affine-finite-morphism-affine-Noetherian} により
$W$ はスキームである。これで証明が完了する。
\end{remark}

\section{スキーム的軌跡と体の拡大}
\label{spaces-over-fields-section-schematic-and-field-extension}

\noindent
体 $k$ 上の表現可能でない代数空間が、$k$ の拡大への基底変換後に
表現可能（すなわちスキーム）となることがある。Spaces, Example
\ref{spaces-example-non-representable-descent} を参照せよ。
本節ではこの問題を扱う。

\begin{lemma}
\label{spaces-over-fields-lemma-scheme-after-purely-inseparable-base-change}
$k$ を体とし、$X$ を $k$ 上の代数空間とする。純非分離体拡大 $k'/k$ で
$X_{k'}$ がスキームとなるものが存在するならば、$X$ はスキームである。
\end{lemma}

\begin{proof}
射 $X_{k'} \to X$ は整、全射、かつ普遍的単射である。したがって本補題は
Limits of Spaces, Lemma
\ref{spaces-limits-lemma-integral-universally-bijective-scheme}
から従う。
\end{proof}

\begin{lemma}
\label{spaces-over-fields-lemma-when-scheme-after-base-change}
$k$ を体、$\overline{k}$ をその代数閉包とする。$X$ を $k$ 上の
準分離代数空間とする。
\begin{enumerate}
\item 体拡大 $K/k$ で $X_K$ がスキームとなるものが存在するならば、
$X_{\overline{k}}$ はスキームである。
\item $X$ が準コンパクトであり、体拡大 $K/k$ で $X_K$ がスキームとなる
ものが存在するならば、$X_{k'}$ は、$k'$ を $k$ のある有限分離拡大として
選ぶことによりスキームとなる。
\end{enumerate}
\end{lemma}

\begin{proof}
任意の代数空間はその準コンパクト開部分空間の合併なので、補題の第 1 の主張は
第 2 の主張から従う（詳細の一部は省略する）。そこで $X$ は準コンパクトであり、
拡大 $K/k$ が与えられ、その基底変換 $X_K$ が表現可能であると仮定する。
$K = \bigcup A$ を有限生成 $k$-部分代数 $A$ の余極限として書く。
Limits of Spaces, Lemma \ref{spaces-limits-lemma-limit-is-scheme} により、
$X_A$ はある $A$ に対してスキームである。極大イデアル
$\mathfrak m \subset A$ を選ぶ。Hilbert の零点定理
（Algebra, Theorem \ref{algebra-theorem-nullstellensatz}）により、
剰余体 $k' = A/\mathfrak m$ は $k$ の有限拡大である。したがって
$X_{k'}$ はスキームである。$k' \supset k$ が分離的でなければ、
Fields, Lemma \ref{fields-lemma-separable-first} で得られる中間拡大を
$k'/k''/k$ とする。$k'/k''$ は純非分離なので、Lemma
\ref{spaces-over-fields-lemma-scheme-after-purely-inseparable-base-change} により代数空間
$X_{k''}$ はスキームである。$k''|k$ は分離的なので証明は完了する。
% P09-ERR-0271
\end{proof}

\begin{lemma}
\label{spaces-over-fields-lemma-base-change-by-Galois}
$k'/k$ を Galois 群 $G$ をもつ有限 Galois 拡大とする。
$X$ を $k$ 上の代数空間とする。このとき $G$ は代数空間 $X_{k'}$ に
自由に作用し、Properties of Spaces, Lemma
\ref{spaces-properties-lemma-quotient} の意味で $X = X_{k'}/G$ である。
\end{lemma}

\begin{proof}
省略する。ヒント：まず $\Spec(k) = \Spec(k')/G$ を示し、次に商をとる操作と
基底変換との両立性を用いる。
\end{proof}

\begin{lemma}
\label{spaces-over-fields-lemma-when-quotient-scheme-at-point}
$S$ をスキームとする。$X$ を $S$ 上の代数空間とし、$G$ を $X$ に自由に
作用する有限群とする。Properties of Spaces, Lemma
\ref{spaces-properties-lemma-quotient} に従い $Y = X/G$ とおく。
$y \in |Y|$ に対し、次は同値である。
\begin{enumerate}
\item $y$ は $Y$ のスキーム的軌跡に属する。
\item アフィン開部分集合 $U \subset X$ で、$y$ の逆像を含むものが存在する。
\end{enumerate}
\end{lemma}

\begin{proof}
Properties of Spaces, Lemma \ref{spaces-properties-lemma-quotient} における
$Y = X/G$ の構成から、射 $X \to Y$ は全射かつ \'etale である。
もちろん $X \times_Y X = X \times G$ であるから、射 $X \to Y$ は実際には
有限 \'etale である。またこれは全射である。したがって本補題は
Decent Spaces, Lemma
\ref{decent-spaces-lemma-when-quotient-scheme-at-point} から従う。
\end{proof}

\begin{lemma}
\label{spaces-over-fields-lemma-scheme-after-purely-transcendental-base-change}
$k$ を体とし、$X$ を $k$ 上の準分離代数空間とする。純超越体拡大 $K/k$ で
$X_K$ がスキームとなるものが存在するならば、$X$ はスキームである。
\end{lemma}

\begin{proof}
任意の代数空間はその準コンパクト開部分空間の合併なので、$X$ は
準コンパクトであると仮定してよい（詳細の一部は省略する）。
拡大 $K/k$ に関する仮定は、$K$ が $k$ 上の多項式環（変数は無限個でもよい）
の分数体であることを意味することを思い出そう（Fields, Definition
\ref{fields-definition-transcendence}）。したがって $K = \bigcup A$ は
部分代数の合併であり、その各々は $k$ 上の有限変数多項式代数の局所化である。
Limits of Spaces, Lemma \ref{spaces-limits-lemma-limit-is-scheme} により、
$X_A$ がスキームとなるような $A$ を選べる。次に
$$
A = k[x_1, \ldots, x_n][1/f]
$$
と書く。ここで $f \in k[x_1, \ldots, x_n]$ は非零である。

\medskip\noindent
$k$ が無限体なら、次のように証明を終えられる：
$a_1, \ldots, a_n \in k$ を
$f(a_1, \ldots, a_n) \not = 0$ となるように選ぶ。
このとき $(a_1, \ldots, a_n)$ は $k$-代数写像 $A \to k$ を定め、
これは $x_i$ を $a_i$ に、$1/f$ を $1/f(a_1, \ldots, a_n)$ に写す。
したがって基底変換
$X_A \times_{\Spec(A)} \Spec(k) \cong X$ は、望むとおりスキームである。
% P09-ERR-0272

\medskip\noindent
この段落では $k$ が有限体の場合に証明を終える。この場合、
$X = \lim X_i$ と書く。ここで $X_i$ は $k$ 上有限表示であり、遷移射は
アフィンである（Limits of Spaces, Lemma
\ref{spaces-limits-lemma-relative-approximation}）。
Limits of Spaces, Lemma \ref{spaces-limits-lemma-limit-is-scheme} を用いると、
$X_{i, A}$ はある $i$ に対してスキームである。よって
$X \to \Spec(k)$ は有限表示であると仮定してよい。
$x \in |X|$ を閉点とする。$x$ は閉埋め込み
$\Spec(\kappa) \to X$ によって表せる
（Decent Spaces, Lemma
\ref{decent-spaces-lemma-decent-space-closed-point}）。
すると $\Spec(\kappa) \to \Spec(k)$ は有限型であるから、$\kappa$ は
$k$ の有限拡大である（Hilbert の零点定理による。Algebra, Theorem
\ref{algebra-theorem-nullstellensatz} を参照；詳細の一部は省略する）。
$[\kappa : k] = d$ とする。整数 $n \gg 0$ を $d$ と互いに素となるように選び、
$k'/k$ を次数 $n$ の拡大とする。このとき $k'/k$ は Galois 拡大であり、
$G = \text{Aut}(k'/k)$ は位数 $n$ の巡回群である。$n$ が十分大きければ、
上と同じ理由により $k$-代数準同型 $A \to k'$ が存在する。
% P09-ERR-0273
すると $X_{k'}$ はスキームであり、Lemma
\ref{spaces-over-fields-lemma-base-change-by-Galois} により $X = X_{k'}/G$ である。
一方、$n$ と $d$ は互いに素なので、
$$
\Spec(\kappa) \times_{X} X_{k'} =
\Spec(\kappa) \times_{\Spec(k)} \Spec(k') =
\Spec(\kappa \otimes_k k')
$$
は体のスペクトルである。言い換えれば、$X_{k'} \to X$ の $x$ 上のファイバーは
一点からなる。したがって Lemma \ref{spaces-over-fields-lemma-when-quotient-scheme-at-point} により、
$x$ は望むとおり $X$ のスキーム的軌跡に属する。
\end{proof}

\begin{remark}
\label{spaces-over-fields-remark-when-does-the-argument-work}
$k$ を有限体とし、$K/k$ を幾何学的に既約な体拡大とする。このとき $K$ は、
幾何学的に既約な有限型 $k$-代数 $A$ の極限である。$A$ が与えられたとき、
Lang と Weil の評価 \cite{LW} により、$n \gg 0$ ならば
$k$-代数準同型 $A \to k'$ が存在し、$k'/k$ の次数は $n$ である。
% P09-ERR-0274
Lemma \ref{spaces-over-fields-lemma-scheme-after-purely-transcendental-base-change} の証明で
与えた議論を解析すると、$X$ が $k$ 上の準分離代数空間であり $X_K$ が
スキームならば、$X$ はスキームであることが分かる。この結果が必要になれば、
ここで正確に定式化して証明することにする。
% P09-ERR-0275
\end{remark}

\begin{lemma}
\label{spaces-over-fields-lemma-scheme-over-algebraic-closure-enough-affines}
$k$ を体、$\overline{k}$ をその代数閉包とする。$X$ を $k$ 上の
代数空間とし、次を仮定する。
\begin{enumerate}
\item $X$ は decent であり、$k$ 上局所有限型である。
\item $X_{\overline{k}}$ はスキームである。
\item $\overline{k}$-有理点からなる $X_{\overline{k}}$ の任意の有限集合は、
あるアフィン開部分集合に含まれる。
% P09-ERR-0276
\end{enumerate}
このとき $X$ はスキームである。
\end{lemma}

\begin{proof}
$K/k$ が拡大ならば、基底変換 $X_K$ は decent であり
（Decent Spaces, Lemma
\ref{decent-spaces-lemma-representable-named-properties}）、また $K$ 上局所有限型である
（Morphisms of Spaces, Lemma
\ref{spaces-morphisms-lemma-base-change-finite-type}）。
Lemma \ref{spaces-over-fields-lemma-scheme-after-purely-inseparable-base-change} により、
$X$ が $k$ の完全化への基底変換後にスキームとなることを示せば十分である。
したがって $k$ は完全体であると仮定してよい（この段階は厳密には不要だが、
残りの議論を考えやすくする）。$X$ を準コンパクト開部分空間で被覆することにより、
$X$ が準コンパクトの場合を証明すれば十分である（細部を一つ省略する）。
この場合 $|X|$ は sober な位相空間である
（Decent Spaces, Proposition
\ref{decent-spaces-proposition-reasonable-sober}）。
したがって $|X|$ のすべての閉点が $X$ のスキーム的軌跡に含まれることを示せば
十分である（Properties of Spaces, Lemma
\ref{spaces-properties-lemma-subscheme} および Topology, Lemma
\ref{topology-lemma-quasi-compact-closed-point} を用いよ）。

\medskip\noindent
$x \in |X|$ を閉点とする。Decent Spaces, Lemma
\ref{decent-spaces-lemma-decent-space-closed-point} により、閉埋め込み
$\Spec(l) \to X$ で $x$ を表すものを見つけられる。このとき
$\Spec(l) \to \Spec(k)$ は有限型であり
（Morphisms of Spaces, Lemma
\ref{spaces-morphisms-lemma-composition-finite-type}）、Hilbert の零点定理
（Algebra, Theorem \ref{algebra-theorem-nullstellensatz}）により
$l$ は $k$ の有限拡大である。$k$ は完全体なので、これは分離拡大である。
したがってスキーム
$$
\Spec(l) \times_X X_{\overline{k}} =
\Spec(l) \times_{\Spec(k)} \Spec(\overline{k}) =
\Spec(l \otimes_k \overline{k})
$$
は有限個の $\overline{k}$-有理点の非交和である。仮定 (3) により、
これらの点を含むアフィン開部分集合 $W \subset X_{\overline{k}}$ が存在する。

\medskip\noindent
Lemma \ref{spaces-over-fields-lemma-when-scheme-after-base-change} により、$X_{k'}$ はある有限拡大
$k'/k$ に対してスキームである。$k'$ を拡大することにより、
アフィン開部分集合 $U' \subset X_{k'}$ が存在し、その
$\overline{k}$ への基底変換が $W$ となると仮定してよい
（$X_{\overline{k}}$ はスキーム $X_{k''}$ の極限であり、ここで
$k' \subset k'' \subset \overline{k}$ は有限拡大であること、および
Limits, Lemmas \ref{limits-lemma-descend-opens} and
\ref{limits-lemma-limit-affine} を用いる）。
$k'/k$ は Galois 拡大であると仮定してよい
（正規閉包をとる。Fields, Lemma
\ref{fields-lemma-normal-closure} を参照し、$k$ が完全体であることを用いる）。
% P09-ERR-0277
$G = \text{Gal}(k'/k)$ とおく。構成により、$G$-不変閉部分スキーム
$\Spec(l) \times_X X_{k'}$ は $U'$ に含まれる。したがって Lemmas
\ref{spaces-over-fields-lemma-base-change-by-Galois} および
\ref{spaces-over-fields-lemma-when-quotient-scheme-at-point} により、$x$ はスキーム的軌跡に属する。
\end{proof}

\noindent
次の二つの補題は別の場所に置くべきである。次の補題を Decent Spaces, Lemma
\ref{decent-spaces-lemma-conditions-on-space-over-field} と比較されたい。
% P09-ERR-0278

\begin{lemma}
\label{spaces-over-fields-lemma-locally-quasi-finite-over-field}
$k$ を体とし、$X$ を $k$ 上の代数空間とする。次は同値である。
\begin{enumerate}
\item $X$ は $k$ 上局所準有限である。
\item $X$ は $k$ 上局所有限型であり、次元 $0$ である。
\item $X$ はスキームであり、$k$ 上局所準有限である。
\item $X$ はスキームであり、$k$ 上局所有限型で、次元 $0$ である。
\item $X$ は Artin 局所 $k$-代数 $A$ のスペクトルの非交和であり、
それらは $k$ 上で $\dim_k(A) < \infty$ を満たす。
\end{enumerate}
\end{lemma}

\begin{proof}
体上で考えているので、$X/k$ の相対次元は $X$ の次元と同じである。
% P09-ERR-0279
したがって Morphisms of Spaces, Lemma
\ref{spaces-morphisms-lemma-locally-quasi-finite-rel-dimension-0}
により、(1) と (2) は同値である。ゆえに Lemma
\ref{spaces-over-fields-lemma-locally-finite-type-dim-zero}（および自明な含意）から、
(1) -- (4) は同値である。最後に Varieties, Lemma
\ref{varieties-lemma-algebraic-scheme-dim-0} により、(1) -- (4) は
(5) と同値である。
% P09-ERR-0280
\end{proof}

\begin{lemma}
\label{spaces-over-fields-lemma-mono-towards-locally-quasi-finite-over-field}
$k$ を体とする。$f : X \to Y$ を $k$ 上の代数空間の単射とする。
$Y$ が $k$ 上局所準有限ならば、$X$ も同様である。
\end{lemma}

\begin{proof}
$Y$ は $k$ 上局所準有限であると仮定する。Lemma
\ref{spaces-over-fields-lemma-locally-quasi-finite-over-field} により、
$Y = \coprod \Spec(A_i)$ であり、各 $A_i$ は $k$ 上有限な Artin 局所環である。
Decent Spaces, Lemma
\ref{decent-spaces-lemma-monomorphism-toward-disjoint-union-dim-0-rings}
により、$X$ はスキームである。$X_i = f^{-1}(\Spec(A_i))$ を考える。
$X_i$ の点は一つまたは零個である。$X_i$ の点が零個なら証明すべきことはない。
$X_i$ の点が一つならば、$X_i = \Spec(B_i)$ である。ここで $B_i$ は
零次元局所環であり、$A_i \to B_i$ は環のエピ射である。特に
$A_i/\mathfrak m_{A_i} = B_i/\mathfrak m_{A_i}B_i$ であり、
$A_i \to B_i$ は中山の補題、Algebra, Lemma
\ref{algebra-lemma-NAK} により全射である
（$\mathfrak m_{A_i}$ は冪零イデアルである！）。したがって
$B_i$ は有限局所 $k$-代数であり、
Lemma \ref{spaces-over-fields-lemma-locally-quasi-finite-over-field} により
$X \to \Spec(k)$ は局所準有限である。
\end{proof}











\section{幾何学的被約代数空間}
\label{spaces-over-fields-section-geometrically-reduced}

\noindent
$X$ が体上の被約代数空間であっても、基礎体を拡大すると $X$ が非被約となる
ことがある。幾何学的被約代数空間ではこの現象は起こらない。

\begin{definition}
\label{spaces-over-fields-definition-geometrically-reduced}
$k$ を体とする。$X$ を $k$ 上の代数空間とする。
\begin{enumerate}
\item $x \in |X|$ を点とする。$X$ が
{\it $x$ で幾何学的被約}であるとは、$\mathcal{O}_{X, \overline{x}}$ が
$k$ 上幾何学的被約であることをいう。
\item $X$ が $k$ 上{\it 幾何学的被約}であるとは、$X$ が
$X$ のすべての点で幾何学的被約であることをいう。
\end{enumerate}
\end{definition}

\noindent
$X$ が $x$ で幾何学的被約ならば、$X$ の $x$ における局所環は被約である
（Properties of Spaces, Lemma
\ref{spaces-properties-lemma-reduced-local-ring}）。同様に、$X$ が
$k$ 上幾何学的被約ならば、$X$ は被約である
（Properties of Spaces, Lemma \ref{spaces-properties-lemma-reduced-space}）。
特に次の補題は、この定義が体上のスキームに対する対応する性質と矛盾しないことを示す。

\begin{lemma}
\label{spaces-over-fields-lemma-geometrically-reduced-at-point}
$k$ を体とし、$X$ を $k$ 上の代数空間とする。$x \in |X|$ とする。
次は同値である。
\begin{enumerate}
\item $X$ は $x$ で幾何学的被約である。
\item ある \'etale 近傍 $(U, u) \to (X, x)$ で、$U$ がスキームであり、
$U$ が $u$ で幾何学的被約となるものが存在する。
\item 任意の \'etale 近傍 $(U, u) \to (X, x)$ で $U$ がスキームであるものに対し、
$U$ は $u$ で幾何学的被約である。
\end{enumerate}
\end{lemma}

\begin{proof}
局所環 $\mathcal{O}_{X, \overline{x}}$ は $\mathcal{O}_{U, u}$ の
厳密 Hensel 化であることを思い出そう。Properties of Spaces, Lemma
\ref{spaces-properties-lemma-describe-etale-local-ring} を参照せよ。
Varieties, Lemma \ref{varieties-lemma-geometrically-reduced-at-point} により、
$U$ が $u$ で幾何学的被約であることは、$\mathcal{O}_{U, u}$ が
$k$ 上幾何学的被約であることと同値である。したがって次を示す必要がある：
$A$ が局所 $k$-代数ならば、$A$ が $k$ 上幾何学的被約であることと
$A^{sh}$ が $k$ 上幾何学的被約であることは同値である。
これを幾何学的被約代数の定義
（Algebra, Definition \ref{algebra-definition-geometrically-reduced}）
を用いて確認する。$K/k$ を体拡大とする。$A \to A^{sh}$ は忠実平坦なので
（More on Algebra, Lemma
\ref{more-algebra-lemma-dumb-properties-henselization}）、
$A \otimes_k K \to A^{sh} \otimes_k K$ は忠実平坦である
（Algebra, Lemma \ref{algebra-lemma-flat-base-change}）。
したがって $A^{sh} \otimes_k K$ が被約ならば、
Algebra, Lemma \ref{algebra-lemma-descent-reduced} により
$A \otimes_k K$ も被約である。逆に、$A^{sh}$ は \'etale
$A$-代数の余極限であることを思い出そう。Algebra, Lemma
% P09-ERR-0281
\ref{algebra-lemma-strict-henselization} を参照せよ。したがって
$A^{sh} \otimes_k K$ は \'etale $A \otimes_k K$-代数のフィルター余極限である。
Algebra, Lemma \ref{algebra-lemma-reduced-goes-up} により結論を得る。
\end{proof}

\begin{lemma}
\label{spaces-over-fields-lemma-geometrically-reduced}
$k$ を体とし、$X$ を $k$ 上の代数空間とする。次は同値である。
\begin{enumerate}
\item $X$ は幾何学的被約である。
\item ある全射 \'etale 射 $U \to X$ で、$U$ がスキームかつ
$U$ が幾何学的被約となるものが存在する。
\item 任意の \'etale 射 $U \to X$ で、$U$ がスキームとなるものに対し、
$U$ は幾何学的被約である。
\end{enumerate}
\end{lemma}

\begin{proof}
定義と Lemma \ref{spaces-over-fields-lemma-geometrically-reduced-at-point} から直ちに従う。
\end{proof}

\noindent
この概念は標数零では興味をもたない。

\begin{lemma}
\label{spaces-over-fields-lemma-perfect-reduced}
$X$ を完全体 $k$ 上の代数空間とする（たとえば $k$ の標数が零である場合）。
\begin{enumerate}
\item $x \in |X|$ に対し、$\mathcal{O}_{X, \overline{x}}$ が被約ならば、
$X$ は $x$ で幾何学的被約である。
\item $X$ が被約ならば、$X$ は $k$ 上幾何学的被約である。
\end{enumerate}
\end{lemma}

\begin{proof}
第 1 の主張は Algebra, Lemma
\ref{algebra-lemma-separable-extension-preserves-reducedness} と完全体の定義
（Algebra, Definition \ref{algebra-definition-perfect}）から従う。
第 2 の主張は第 1 の主張から従う。
\end{proof}

\begin{lemma}
\label{spaces-over-fields-lemma-geometrically-reduced-positive-characteristic}
$k$ を標数 $p > 0$ の体とし、$X$ を $k$ 上の代数空間とする。
次は同値である。
\begin{enumerate}
\item $X$ は $k$ 上幾何学的被約である。
\item $X_{k'}$ は任意の体拡大 $k'/k$ に対して被約である。
\item $X_{k'}$ は任意の有限純非分離体拡大 $k'/k$ に対して被約である。
\item $X_{k^{1/p}}$ は被約である。
\item $X_{k^{perf}}$ は被約である。
\item $X_{\bar k}$ は被約である。
\end{enumerate}
\end{lemma}

\begin{proof}
全射 \'etale 射 $U \to X$ で $U$ がスキームとなるものを選ぶ。
Lemma \ref{spaces-over-fields-lemma-geometrically-reduced} を介して、本補題は $U$ の $k$ 上の
対応する結果から従う。Varieties, Lemma
\ref{varieties-lemma-geometrically-reduced} を参照せよ。
\end{proof}

\begin{lemma}
\label{spaces-over-fields-lemma-geometrically-reduced-upstairs}
$k$ を体とし、$X$ を $k$ 上の代数空間とする。$k'/k$ を体拡大とする。
$x \in |X|$ を点とし、$x' \in |X_{k'}|$ を $x$ の上にある点とする。
次は同値である。
\begin{enumerate}
\item $X$ は $x$ で幾何学的被約である。
\item $X_{k'}$ は $x'$ で幾何学的被約である。
\end{enumerate}
特に、$X$ が $k$ 上幾何学的被約であることと、$X_{k'}$ が $k'$ 上
幾何学的被約であることは同値である。
\end{lemma}

\begin{proof}
\'etale 射 $U \to X$ で $U$ がスキームとなるものを選び、さらに点
$u \in U$ で $x \in |X|$ に写るものを選ぶ。Properties of Spaces, Lemma
\ref{spaces-properties-lemma-points-cartesian} により、点
$u' \in U_{k'} = U \times_X X_{k'}$ で $u$ と $x'$ の両方へ写るものを選べる。Lemma
\ref{spaces-over-fields-lemma-geometrically-reduced-at-point} により、本補題は $U, u, u'$ に対する
補題、すなわち Varieties, Lemma
\ref{varieties-lemma-geometrically-reduced-upstairs} から従う。
\end{proof}

\begin{lemma}
\label{spaces-over-fields-lemma-geometrically-reduced-etale-local}
$k$ を体とする。$f : X \to Y$ を $k$ 上の代数空間の射とする。
$x \in |X|$ を点とし、その像を $y \in |Y|$ とする。
\begin{enumerate}
\item $f$ が $x$ で \'etale ならば、
$X$ が $x$ で幾何学的被約 $\Leftrightarrow$
$Y$ が $y$ で幾何学的被約である。
\item $f$ が全射 \'etale ならば、
$X$ が幾何学的被約 $\Leftrightarrow$
$Y$ が幾何学的被約である。
\end{enumerate}
\end{lemma}

\begin{proof}
(1) は明らかである。実際、
$\mathcal{O}_{X, \overline{x}} = \mathcal{O}_{Y, \overline{y}}$
である。ただし $f$ は $x$ で \'etale である。(2) は (1) から直ちに従う。
\end{proof}







\section{幾何学的連結代数空間}
\label{spaces-over-fields-section-geometrically-connected}

\noindent
$X$ が体上の連結代数空間であっても、基礎体を拡大すると $X$ が非連結となる
ことがある。幾何学的連結代数空間ではこの現象は起こらない。

\begin{definition}
\label{spaces-over-fields-definition-geometrically-connected}
$X$ を体 $k$ 上の代数空間とする。$X$ が $k$ 上
{\it 幾何学的連結}であるとは、基底変換 $X_{k'}$ が、
$k'$ を $k$ の任意の体拡大として連結であることをいう。
\end{definition}

\noindent
規約により連結位相空間は空でない。したがって当然、幾何学的連結代数空間は
空でない。

\begin{lemma}
\label{spaces-over-fields-lemma-geometrically-connected-check-after-extension}
$X$ を体 $k$ 上の代数空間とし、$k'/k$ を体拡大とする。このとき
$X$ が $k$ 上幾何学的連結であることと、$X_{k'}$ が $k'$ 上
幾何学的連結であることは同値である。
\end{lemma}

\begin{proof}
$X$ が $k$ 上幾何学的連結ならば、$X_{k'}$ が $k'$ 上幾何学的連結であることは
明らかである。逆を示すため、任意の体拡大 $k''/k$ に対し、共通拡大体
$k'''/k$ が $k'/k$ と $k''/k$ に対して存在する。Fields, Lemma
\ref{fields-lemma-common-extension-field} を参照せよ。射
$X_{k'''} \to X_{k''}$ は全射なので（体のスペクトル間の全射の基底変換である）、
$X_{k'''}$ の連結性から $X_{k''}$ の連結性が従う。したがって
$X_{k'}$ が $k'$ 上幾何学的連結ならば、$X$ は $k$ 上幾何学的連結である。
\end{proof}

\begin{lemma}
\label{spaces-over-fields-lemma-bijection-connected-components}
$k$ を体とし、$X$, $Y$ を $k$ 上の代数空間とする。
$X$ は $k$ 上幾何学的連結であると仮定する。このとき射影
$$
p : X \times_k Y \longrightarrow Y
$$
は連結成分の間の全単射を誘導する。
\end{lemma}

\begin{proof}
$y \in |Y|$ は射 $\Spec(K) \to Y$ で表され、$K$ は体であるとする。
Properties of Spaces, Lemma \ref{spaces-properties-lemma-points-cartesian} により、
$|X \times_k Y| \to |Y|$ の $y$ 上のファイバーは
$|X_K| \to |X \times_k Y|$ の像である。したがって、$X$ が幾何学的連結である
という仮定により、これらのファイバーは連結である。Morphisms of Spaces, Lemma
\ref{spaces-morphisms-lemma-space-over-field-universally-open} により、写像
$|p|$ は開である。よって Topology, Lemma
\ref{topology-lemma-connected-fibres-connected-components} を適用して結論を得る。
\end{proof}

\begin{lemma}
\label{spaces-over-fields-lemma-separably-closed-field-connected-components}
$k'/k$ を体拡大とし、$X$ を $k$ 上の代数空間とする。$k$ は分離代数閉であると
仮定する。このとき射 $X_{k'} \to X$ は連結成分の全単射を誘導する。
特に、$X$ が $k$ 上幾何学的連結であることと $X$ が連結であることは同値である。
\end{lemma}

\begin{proof}
$k$ は分離代数閉なので、$k'$ は $k$ 上幾何学的連結である。Algebra, Lemma
\ref{algebra-lemma-separably-closed-connected-implies-geometric} を参照せよ。
したがって Varieties, Lemma
\ref{varieties-lemma-affine-geometrically-connected} により
$Z = \Spec(k')$ は $k$ 上幾何学的連結である。
$X_{k'} = Z \times_k X$ なので、結論は Lemma
\ref{spaces-over-fields-lemma-bijection-connected-components} の特別な場合である。
\end{proof}

\begin{lemma}
\label{spaces-over-fields-lemma-characterize-geometrically-connected}
$k$ を体とし、$X$ を $k$ 上の代数空間とする。$\overline{k}$ を $k$ の
分離代数閉包とする。このとき $X$ が幾何学的連結であることと、基底変換
$X_{\overline{k}}$ が連結であることは同値である。
\end{lemma}

\begin{proof}
$X_{\overline{k}}$ が連結であると仮定する。$k'/k$ を体拡大とする。
体拡大 $\overline{k}'/\overline{k}$ で、$k'$ が $\overline{k}'$ に
$k$ の拡大として埋め込まれるものが存在する。Lemma
\ref{spaces-over-fields-lemma-separably-closed-field-connected-components} により
$X_{\overline{k}'}$ は連結である。$X_{\overline{k}'} \to X_{k'}$ は全射なので、
$X_{k'}$ は望むとおり連結である。
\end{proof}

\noindent
$k$ を体とし、$\overline{k}/k$ を（無限でもよい）Galois 拡大とする。
たとえば $\overline{k}$ は $k$ の分離代数閉包でもよい。
任意の $\sigma \in \text{Gal}(\overline{k}/k)$ に対し、対応する自己同型
$
\Spec(\sigma) :
\Spec(\overline{k})
\longrightarrow
\Spec(\overline{k})
$ が得られる。ここで
$\Spec(\sigma) \circ \Spec(\tau) = \Spec(\tau \circ \sigma)$ であることに注意せよ。
したがって作用
$$
\text{Gal}(\overline{k}/k)^{opp} \times \Spec(\overline{k})
\longrightarrow
\Spec(\overline{k})
$$
が得られる。これは反対群のスキーム $\Spec(\overline{k})$ 上の作用である。
$X$ を $k$ 上の代数空間とする。定義により
$X_{\overline{k}} =
\Spec(\overline{k}) \times_{\Spec(k)} X$
なので、上の作用は標準的な作用
\begin{equation}
\label{spaces-over-fields-equation-galois-action-base-change-kbar}
\text{Gal}(\overline{k}/k)^{opp} \times X_{\overline{k}}
\longrightarrow
X_{\overline{k}}.
\end{equation}
を誘導する。

\begin{lemma}
\label{spaces-over-fields-lemma-Galois-action-quasi-compact-open}
$k$ を体とし、$X$ を $k$ 上の代数空間とする。$\overline{k}$ を $k$ の
（無限でもよい）Galois 拡大とする。$V \subset X_{\overline{k}}$ を
準コンパクト開部分空間とする。このとき、
\begin{enumerate}
\item 有限中間拡大 $\overline{k}/k'/k$ と準コンパクト開部分空間
$V' \subset X_{k'}$ で $V = (V')_{\overline{k}}$ となるものが存在する。
\item 開部分群 $H \subset \text{Gal}(\overline{k}/k)$ で、
$\sigma(V) = V$ がすべての $\sigma \in H$ に対して成立するものが存在する。
\end{enumerate}
\end{lemma}

\begin{proof}
スキーム $U$ と全射 \'etale 射 $U \to X$ を選ぶ。準コンパクト開部分空間
$W \subset U_{\overline{k}}$ で、その $X_{\overline{k}}$ における像が
$V$ となるものを選ぶ。これは $|U_{\overline{k}}| \to |X_{\overline{k}}|$ が
連続であり、$|U_{\overline{k}}|$ が準コンパクト開集合からなる基底をもつので可能である。
$W \subset U_{\overline{k}}$ に Varieties, Lemma
\ref{varieties-lemma-Galois-action-quasi-compact-open} を適用して補題を得る。
\end{proof}

\begin{lemma}
\label{spaces-over-fields-lemma-closed-fixed-by-Galois}
$k$ を体とし、$\overline{k}/k$ を（無限でもよい）Galois 拡大とする。
$X$ を $k$ 上の代数空間とする。
$\overline{T} \subset |X_{\overline{k}}|$ が次の性質をもつとする。
\begin{enumerate}
\item $\overline{T}$ は $|X_{\overline{k}}|$ の閉部分集合である。
\item すべての $\sigma \in \text{Gal}(\overline{k}/k)$ に対し、
$\sigma(\overline{T}) = \overline{T}$ である。
\end{enumerate}
このとき閉部分集合 $T \subset |X|$ で、その $|X_{k'}|$ における逆像が
$\overline{T}$ となるものが存在する。
% P09-ERR-0282
\end{lemma}

\begin{proof}
$T \subset |X|$ を $\overline{T}$ の像とする。
$|X_{\overline{k}}| \to |X|$ は全射なので、主張は $T$ が閉であり、その逆像が
$\overline{T}$ であることを意味する。スキーム $U$ と全射 \'etale 射
$U \to X$ を選ぶ。スキームの場合
（Varieties, Lemma \ref{varieties-lemma-closed-fixed-by-Galois} を参照）
により、閉部分集合 $T' \subset |U|$ で、その $|U_{\overline{k}}|$ における
逆像が $\overline{T}$ の逆像となるものが存在する。
$|U_{\overline{k}}| \to |X_{\overline{k}}|$ は全射なので、$T'$ は
$T$ の逆像であり、その写像は $|U| \to |X|$ である。$|X|$ の位相の構成により、
これは $T$ が閉であることを意味する。同じ方法で $\overline{T}$ が
$T$ の逆像であることが分かる。
\end{proof}

\begin{lemma}
\label{spaces-over-fields-lemma-characterize-geometrically-disconnected}
$k$ を体とし、$X$ を $k$ 上の代数空間とする。次は同値である。
\begin{enumerate}
\item $X$ は幾何学的連結である。
\item 任意の有限分離体拡大 $k'/k$ に対し、代数空間 $X_{k'}$ は連結である。
\end{enumerate}
\end{lemma}

\begin{proof}
この証明は Varieties, Lemma
\ref{varieties-lemma-characterize-geometrically-disconnected} の証明と同じである。
ただし、Varieties, Lemma
\ref{varieties-lemma-characterize-geometrically-connected} を Lemma
\ref{spaces-over-fields-lemma-characterize-geometrically-connected} に置き換え、
Varieties, Lemma
\ref{varieties-lemma-Galois-action-quasi-compact-open} を Lemma
\ref{spaces-over-fields-lemma-Galois-action-quasi-compact-open} に置き換え、さらに
Varieties, Lemma
\ref{varieties-lemma-closed-fixed-by-Galois} を Lemma
\ref{spaces-over-fields-lemma-closed-fixed-by-Galois} に置き換える。読者には、こちらではなく
そちらの証明を読むことを強く勧める。
% P09-ERR-0283

\medskip\noindent
定義から直ちに (1) は (2) を含意する。$X$ が幾何学的連結でないと仮定する。
$k \subset \overline{k}$ を $k$ の分離代数閉包とする。Lemma
\ref{spaces-over-fields-lemma-characterize-geometrically-connected} により
$X_{\overline{k}}$ は非連結である。そこで
$X_{\overline{k}} = \overline{U} \amalg \overline{V}$ と書く。ここで
$\overline{U}$ と $\overline{V}$ は $X_{\overline{k}}$ の空でない、
開かつ閉な代数部分空間である。

\medskip\noindent
$W \subset X$ を任意の準コンパクト開部分空間とする。このとき
$W_{\overline{k}} \cap \overline{U}$ と
$W_{\overline{k}} \cap \overline{V}$ は $W_{\overline{k}}$ の
開かつ閉な部分空間である。特に
$W_{\overline{k}} \cap \overline{U}$ と
$W_{\overline{k}} \cap \overline{V}$ は準コンパクトであり、Lemma
\ref{spaces-over-fields-lemma-Galois-action-quasi-compact-open} により、
$W_{\overline{k}} \cap \overline{U}$ と
$W_{\overline{k}} \cap \overline{V}$ はともに有限中間拡大上で定義され、
$\text{Gal}(\overline{k}/k)$ のある開部分群の下で不変である。
以下ではこれを改めて断ることなく用いる。

\medskip\noindent
準コンパクト開部分空間 $W_0 \subset X$ を、
$W_{0, \overline{k}} \cap \overline{U}$ と
$W_{0, \overline{k}} \cap \overline{V}$ の両方が空でないように選ぶ。
有限中間拡大 $\overline{k}/k'/k$ と、開かつ閉な部分集合への分解
$W_{0, k'} = U_0' \amalg V_0'$ を、次を満たすように選ぶ：
$W_{0, \overline{k}} \cap \overline{U} = (U'_0)_{\overline{k}}$ および
$W_{0, \overline{k}} \cap \overline{V} = (V'_0)_{\overline{k}}$。
$H = \text{Gal}(\overline{k}/k') \subset \text{Gal}(\overline{k}/k)$ とおく。
特に
$\sigma(W_{0, \overline{k}} \cap \overline{U}) =
W_{0, \overline{k}} \cap \overline{U}$ であり、$\overline{V}$ についても同様である。

\medskip\noindent
上のように $W_0$, $k'$ を選んだ後、任意の準コンパクト開部分空間
$W \subset X$ に対し
$$
U_W =
\bigcap\nolimits_{\sigma \in H} \sigma(W_{\overline{k}} \cap \overline{U}),
\quad
V_W =
\bigcup\nolimits_{\sigma \in H} \sigma(W_{\overline{k}} \cap \overline{V}).
$$
とおく。さて、$W_{\overline{k}} \cap \overline{U}$ と
$W_{\overline{k}} \cap \overline{V}$ は $\text{Gal}(\overline{k}/k)$ の
ある開部分群によって固定されるので、上の合併と共通部分は有限である。
したがって $U_W$ と $V_W$ はともに開かつ閉な部分空間である。また構成により
$W_{\bar k} = U_W \amalg V_W$ である。

\medskip\noindent
$W \subset W' \subset X$ が準コンパクト開部分空間ならば
$W_{\overline{k}} \cap U_{W'} = U_W$ および
$W_{\overline{k}} \cap V_{W'} = V_W$ であると主張する。検証は省略する。
したがって $U = \bigcup_{W \subset X} U_W$ および
$V = \bigcup_{W \subset X} V_W$ と定めると、
$X_{\overline{k}} = U \amalg V$ は開かつ閉な部分集合の非交和として得られる。
% P09-ERR-0284
$V$ は（局所的に）合併をとって構成されているので、空でないことは明らかである。
一方、$U$ は構成により $W_0 \cap \overline{U}$ を含むので空でない。
% P09-ERR-0285
最後に、$U, V \subset X_{\bar k}$ は構成により閉かつ $H$-不変である。
したがって Lemma \ref{spaces-over-fields-lemma-closed-fixed-by-Galois} により
$U = (U')_{\bar k}$ および $V = (V')_{\bar k}$ である。ここで
$U', V' \subset X_{k'}$ はある閉部分集合である。
明らかに $X_{k'} = U' \amalg V'$ であり、$X_{k'}$ は望むとおり非連結である。
\end{proof}







\section{幾何学的既約代数空間}
\label{spaces-over-fields-section-geometrically-irreducible}

\noindent
Spaces, Example \ref{spaces-example-infinite-product} は、完全な一般性の下で
既約代数空間を考えない方がよいことを示している
\footnote{正確に言えば、「代数空間 $X$ は既約である」と言うとき、おそらく
「位相空間 $|X|$ は既約である」と言うつもりである。}。
decent な（たとえば準分離な）代数空間では、この種の破綻は起こらない。
そこで以下の定義は、代数空間が decent であるという仮定の下でのみ行う。

\begin{definition}
\label{spaces-over-fields-definition-geometrically-irreducible}
$k$ を体とする。$X$ を $k$ 上の decent な代数空間とする。
$X$ が{\it 幾何学的既約}であるとは、位相空間 $|X_{k'}|$ が
任意の体拡大 $k'$（$k$ の拡大）に対して既約であることをいう
\footnote{既約空間は空でない。}。
\end{definition}

\noindent
$X_{k'}$ は decent な代数空間であることに注意せよ
（Decent Spaces, Lemma
\ref{decent-spaces-lemma-representable-named-properties}）。
したがって位相空間 $|X_{k'}|$ は sober である。
Decent Spaces, Proposition
\ref{decent-spaces-proposition-reasonable-sober}。
% P09-ERR-0286





\section{幾何学的整代数空間}
\label{spaces-over-fields-section-geometrically-integral}

\noindent
整代数空間は定義により decent であることを思い出そう。Section
\ref{spaces-over-fields-section-integral-spaces} を参照せよ。

\begin{definition}
\label{spaces-over-fields-definition-geometrically-integral}
$X$ を体 $k$ 上の代数空間とする。$X$ が $k$ 上{\it 幾何学的整}であるとは、
代数空間 $X_{k'}$ がすべての体拡大 $k'$（$k$ の拡大）に対して整であることをいう
（Definition \ref{spaces-over-fields-definition-integral-algebraic-space}）。
\end{definition}

\noindent
特に $X$ は decent な代数空間である。これは幾何学的被約性および
幾何学的既約性と次のように関係する。

\begin{lemma}
\label{spaces-over-fields-lemma-geometrically-integral}
$k$ を体とし、$X$ を $k$ 上の decent な代数空間とする。
$X$ が $k$ 上幾何学的整であることと、$X$ が $k$ 上幾何学的被約かつ
幾何学的既約であることは同値である。
\end{lemma}

\begin{proof}
これは定義から直ちに従う。実際、decency の下では、ここでの整性の概念は
被約かつ既約であることと同値である。
\end{proof}

\begin{lemma}
\label{spaces-over-fields-lemma-proper-geometrically-reduced-global-sections}
$k$ を体とし、$X$ を $k$ 上の固有代数空間とする。
\begin{enumerate}
\item $A = H^0(X, \mathcal{O}_X)$ は有限次元 $k$-代数である。
\item $A = \prod_{i = 1, \ldots, n} A_i$ は Artin 局所 $k$-代数の積であり、
$|X|$ の各連結成分に一つの因子が対応する。
\item $X$ が被約ならば、$A = \prod_{i = 1, \ldots, n} k_i$ は体の積であり、
各体は $k$ の有限拡大である。
\item $X$ が幾何学的被約ならば、$k_i$ は $k$ 上有限分離的である。
\item $X$ が幾何学的連結ならば、$A$ は $k$ 上幾何学的既約である。
\item $X$ が幾何学的既約ならば、$A$ は $k$ 上幾何学的既約である。
\item $X$ が幾何学的被約かつ連結ならば、$A = k$ である。
% P09-ERR-0288
\item $X$ が幾何学的整ならば、$A = k$ である。
\end{enumerate}
\end{lemma}

\begin{proof}
Cohomology of Spaces, Lemma
\ref{spaces-cohomology-lemma-proper-over-affine-cohomology-finite} により、
$A = H^0(X, \mathcal{O}_X)$ は有限次元 $k$-代数である。これは (1) を証明する。

\medskip\noindent
Algebra, Lemma \ref{algebra-lemma-finite-dimensional-algebra} および
Algebra, Proposition \ref{algebra-proposition-dimension-zero-ring} により、
$A$ は局所環の積である。$X = Y \amalg Z$ で $Y$ と $Z$ が $X$ の
開部分空間ならば、冪等元 $e \in A$ を得る。これは $\mathcal{O}_X$ の、
値 $1$ を $Y$ 上で、値 $0$ を $Z$ 上でとる切断による。逆に
$e \in A$ が冪等元ならば、
$|X|$ の対応する分解を得る。最後に、$|X|$ は Noether 位相空間なので
（Morphisms of Spaces, Lemma
\ref{spaces-morphisms-lemma-finite-presentation-noetherian} および
Properties of Spaces, Lemma
\ref{spaces-properties-lemma-Noetherian-topology} による）、その連結成分は開である。
したがって $|X|$ の連結成分は $1$-対-$1$ に、$A$ の原始冪等元と対応する。
これは (2) を証明する。

\medskip\noindent
$X$ が被約ならば、$A$ は被約である
（Properties of Spaces, Lemma \ref{spaces-properties-lemma-reduced-space}）。
したがって局所環 $A_i = k_i$ は被約であり、それゆえ体である
（たとえば Algebra, Lemma
\ref{algebra-lemma-minimal-prime-reduced-ring} による）。
これは (3) を証明する。

\medskip\noindent
$X$ が幾何学的被約ならば、同じことが
$A \otimes_k \overline{k} =
H^0(X_{\overline{k}}, \mathcal{O}_{X_{\overline{k}}})$
について成り立つ（等式については Cohomology of Spaces, Lemma
\ref{spaces-cohomology-lemma-flat-base-change-cohomology} を参照）。
% P09-ERR-0287
これは $k_i \otimes_k \overline{k}$ が体の積であることを含意し、
したがって $k_i/k$ は分離的である。たとえば Algebra, Lemmas
\ref{algebra-lemma-characterize-separable-field-extensions} および
\ref{algebra-lemma-geometrically-reduced-finite-purely-inseparable-extension}
による。これは (4) を証明する。

\medskip\noindent
$X$ が幾何学的連結ならば、
$A \otimes_k \overline{k} =
H^0(X_{\overline{k}}, \mathcal{O}_{X_{\overline{k}}})$
は (2) により零次元局所環であるから、そのスペクトルは一点をもち、特に既約である。
したがって $A$ は幾何学的既約である。これは (5) を証明する。
もちろん (5) は (6) を含意する。

\medskip\noindent
$X$ が幾何学的被約かつ連結ならば、$A = k_1$ は体であり、拡大 $k_1/k$ は
有限分離かつ幾何学的既約である。しかしこのとき
$k_1 \otimes_k \overline{k}$ は $[k_1 : k]$ 個の $\overline{k}$ の積であるから、
$k_1 = k$ を得る。これは (7) を証明する。もちろん (7) は (8) を含意する。
\end{proof}

\begin{lemma}
\label{spaces-over-fields-lemma-characterize-trivial-pic-integral}
$k$ を体とし、$X$ を $k$ 上の固有整代数空間とする。
$\mathcal{L}$ を可逆 $\mathcal{O}_X$-加群とする。
$H^0(X, \mathcal{L})$ と $H^0(X, \mathcal{L}^{\otimes - 1})$ がともに非零ならば、
$\mathcal{L} \cong \mathcal{O}_X$ である。
\end{lemma}

\begin{proof}
$s \in H^0(X, \mathcal{L})$ と $t \in H^0(X, \mathcal{L}^{\otimes - 1})$ を
非零切断とする。$x \in |X|$ を $s$ の台にある点とする。
アフィン \'etale 近傍 $(U, u) \to (X, x)$ を
$\mathcal{L}|_U \cong \mathcal{O}_U$ となるように選ぶ。このとき $s|_U$ は、
被約な（$X$ が被約なので）スキーム $U$ 上の非零正則函数に対応し、したがって
$U$ のある既約成分の一般点で消えない。
% P09-ERR-0289
Decent Spaces, Lemma
\ref{decent-spaces-lemma-decent-generic-points} により、一般点 $\eta$（$|X|$ のもの）は
$s$ の台にある。$t$ についても同じである。するともちろん $st$ は非零でなければ
ならない。実際、$X$ の $\eta$ における局所環は体である
（前掲の補題により局所環は零次元であり、$X$ が被約なので局所環は被約であり、
Algebra, Lemma \ref{algebra-lemma-minimal-prime-reduced-ring} を用いる）。
一方、Lemma \ref{spaces-over-fields-lemma-proper-geometrically-reduced-global-sections} で
$K = H^0(X, \mathcal{O}_X)$ は体であることを見た。
したがって $st$ は至るところ非零であり、
$s : \mathcal{O}_X \to \mathcal{L}$ は同型である。
\end{proof}







\section{次元}
\label{spaces-over-fields-section-dimension}

\noindent
本節では次元に関する議論を続ける。これまでの内容を列挙すると次のとおりである：
\begin{enumerate}
\item 次元は Properties of Spaces, Section
\ref{spaces-properties-section-dimension} で定義されている。
\item 局所環の次元は Properties of Spaces, Section
\ref{spaces-properties-section-dimension-local-ring} で定義されている。
\item Properties of Spaces, Lemmas
\ref{spaces-properties-lemma-dimension-local-ring} および
\ref{spaces-properties-lemma-dimension-decent-invariant-under-etale} にいくつかの結果がある。
\item 相対次元は Morphisms of Spaces, Section
\ref{spaces-morphisms-section-relative-dimension} で定義されている。
\item ファイバーの次元に関する結果は Morphisms of Spaces, Section
\ref{spaces-morphisms-section-dimension-fibres} にある。
\item 次元公式の弱い形は Morphisms of Spaces, Section
\ref{spaces-morphisms-section-dimension-formula} にある。
\item 滑らかさと次元に関する結果は Morphisms of Spaces, Lemma
\ref{spaces-morphisms-lemma-smoothness-dimension-spaces} にある。
\item decent な空間では次元は $\dim(|X|)$ である。Decent Spaces, Lemma
\ref{decent-spaces-lemma-dimension-decent-space}。
% P09-ERR-0290
\item 準有限写像と次元については Decent Spaces, Lemmas
\ref{decent-spaces-lemma-dimension-local-ring-quasi-finite} および
\ref{decent-spaces-lemma-dimension-quasi-finite}。
\end{enumerate}

More on Morphisms of Spaces, Section
\ref{spaces-more-morphisms-section-dimension} では、有限型射のファイバーにおける
次元の跳躍を論じる。

\begin{lemma}
\label{spaces-over-fields-lemma-integral-dimension}
$S$ をスキームとし、$f : X \to Y$ を代数空間の整射とする。このとき
$\dim(X) \leq \dim(Y)$ である。$f$ が全射ならば
$\dim(X) = \dim(Y)$ である。
\end{lemma}

\begin{proof}
$V \to Y$ を全射 \'etale とし、$V$ がスキームとなるように選ぶ。
このとき $U = X \times_Y V$ はスキームであり、$U \to V$ は整である
（$f$ が全射ならば全射でもある）。Properties of Spaces, Lemma
\ref{spaces-properties-lemma-dimension-decent-invariant-under-etale} により
$\dim(X) = \dim(U)$ および $\dim(Y) = \dim(V)$ である。
したがって結論はスキームの場合、すなわち Morphisms, Lemma
\ref{morphisms-lemma-integral-dimension} から従う。
\end{proof}

\begin{lemma}
\label{spaces-over-fields-lemma-alteration-dimension}
$S$ をスキームとし、$f : X \to Y$ を $S$ 上の代数空間の射とする。
次を仮定する。
\begin{enumerate}
\item $Y$ は局所 Noether である。
\item $X$ と $Y$ は整代数空間である。
\item $f$ は支配的である。
\item $f$ は局所有限型である。
\end{enumerate}
$x \in |X|$ と $y \in |Y|$ が一般点ならば、
$$
\dim(X) \leq \dim(Y) + \text{超越次数 }x/y.
$$
$f$ が固有ならば等号が成り立つ。
\end{lemma}

\begin{proof}
$|X|$ と $|Y|$ は既約 sober 位相空間であることを思い出そう。Definition
\ref{spaces-over-fields-definition-integral-algebraic-space} に続く議論を参照せよ。
したがって $f$ が支配的であるとは、$|f|$ が $x$ を $y$ に写すことを意味する。
さらに $x \in |X|$ は、$X$ の局所環の次元が $0$ となる唯一の点である。
Decent Spaces, Lemma
\ref{decent-spaces-lemma-decent-generic-points} を参照せよ。
Morphisms of Spaces, Lemma
\ref{spaces-morphisms-lemma-dimension-formula-general} により、
$X$ の任意の点 $x' \in |X|$ における局所環の次元は、$Y$ の点
$y' = f(x')$ における局所環の次元と $x/y$ の超越次数との和以下である。
$X$ の次元および $Y$ の次元は、それぞれ $x'$ および $y'$ における局所環の
次元の上限なので（Properties of Spaces, Lemma
\ref{spaces-properties-lemma-dimension}）、不等式を得る。

\medskip\noindent
$f$ は固有であると仮定する。$V \subset Y$ を空でない準コンパクト開部分空間とする。
射 $f^{-1}(V) \to V$ に対して等式を証明できれば、$X \to Y$ に対しても
等式を得る。したがって $X$ と $Y$ は準コンパクトであると仮定してよい。
$X$ は局所 Noether な decent 代数空間として準分離であることに注意せよ。
Decent Spaces, Lemma
\ref{decent-spaces-lemma-locally-Noetherian-decent-quasi-separated} を参照せよ。
したがって有限全射 $Y' \to Y$ で $Y'$ がスキームとなるものを選べる。
Limits of Spaces, Proposition
\ref{spaces-limits-proposition-there-is-a-scheme-finite-over} を参照せよ。
$Y'$ を適当な閉部分スキームで置き換えれば、$Y'$ は整であると仮定してよい。
たとえば、より一般的な Lemma \ref{spaces-over-fields-lemma-alteration-contained-in} を参照せよ。
同じ補題により閉部分空間 $X' \subset X \times_Y Y'$ を、$X'$ が整であり、
$X' \to X$ が有限全射となるように選べる。ここで $X'$ も局所 Noether である
（Morphisms of Spaces, Lemma
\ref{spaces-morphisms-lemma-locally-finite-type-locally-noetherian}）。
Limits of Spaces, Proposition
\ref{spaces-limits-proposition-there-is-a-scheme-finite-over} をもう一度用い、
有限全射 $X'' \to X'$ で $X''$ がスキームとなるものを選ぶ。
前と同様に $X''$ は整であると仮定してよい。図式は
$$
\xymatrix{
X'' \ar[d] \ar[r] & X \ar[d]^f \\
Y' \ar[r] & Y
}
$$
である。Lemma \ref{spaces-over-fields-lemma-integral-dimension} により
$\dim(X'') = \dim(X)$ および $\dim(Y') = \dim(Y)$ である。
$X$ と $Y$ はそれぞれ $x$ と $y$ の、スキームとなる開近傍をもつので、
一般点 $x'' \in X''$ と $y' \in Y'$ はそれぞれ $x$ と $y$ に写る唯一の点であり、
剰余体拡大 $\kappa(x'')/\kappa(x)$ と $\kappa(y')/\kappa(y)$ は有限であることが
直ちに分かる。これは $x''/y'$ の超越次数が $x/y$ の超越次数と同じであることを
含意する。したがって等式はスキームの場合、すなわち Morphisms, Lemma
\ref{morphisms-lemma-alteration-dimension} から従う。
% P09-ERR-0291
\end{proof}






\section{体上滑らかな空間}
\label{spaces-over-fields-section-smooth}

\noindent
本節は Varieties, Section \ref{varieties-section-smooth} の類似である。

\begin{lemma}
\label{spaces-over-fields-lemma-smooth-regular}
$k$ を体とし、$X$ を $k$ 上滑らかな代数空間とする。このとき $X$ は
正則代数空間である。
\end{lemma}

\begin{proof}
スキーム $U$ と全射 \'etale 射 $U \to X$ を選ぶ。射
$U \to \Spec(k)$ は \'etale（したがって滑らかな）射と滑らかな射との合成なので
滑らかである（Morphisms of Spaces, Lemmas
\ref{spaces-morphisms-lemma-etale-smooth} および
\ref{spaces-morphisms-lemma-composition-smooth} を参照）。
したがって Varieties, Lemma \ref{varieties-lemma-smooth-regular} により
$U$ は正則である。Properties of Spaces, Definition
\ref{spaces-properties-definition-type-property} によれば、これは $X$ が
正則であることを意味する。
\end{proof}

\begin{lemma}
\label{spaces-over-fields-lemma-smooth-separable-closed-points-dense}
$k$ を体とし、$X$ を $\Spec(k)$ 上滑らかな代数空間とする。
$x \in |X|$ で射 $\Spec(k') \to X$ の像となり、$k' \supset k$ が有限分離である
もの全体は $|X|$ で稠密である。
% P09-ERR-0292
\end{lemma}

\begin{proof}
スキーム $U$ と全射 \'etale 射 $U \to X$ を選ぶ。射
$U \to \Spec(k)$ は \'etale（したがって滑らかな）射と滑らかな射との合成なので
滑らかである（Morphisms of Spaces, Lemmas
\ref{spaces-morphisms-lemma-etale-smooth} および
\ref{spaces-morphisms-lemma-composition-smooth} を参照）。
したがって Varieties, Lemma
\ref{varieties-lemma-smooth-separable-closed-points-dense} を適用すると、
$U$ の閉点で、その剰余体が $k$ 上有限分離的であるものは稠密である。
これは $|X|$ の位相の定義により本補題を含意する。
\end{proof}








\section{Euler 標数}
\label{spaces-over-fields-section-euler}

\noindent
本節では、体上固有な代数空間上の連接層の Euler 標数について、いくつかの
初等的性質を証明する。

\begin{definition}
\label{spaces-over-fields-definition-euler-characteristic}
$k$ を体とし、$X$ を $k$ 上の固有代数空間とする。
$\mathcal{F}$ を連接 $\mathcal{O}_X$-加群とする。この状況で
{\it $\mathcal{F}$ の Euler 標数}とは整数
% P09-ERR-0293
$$
\chi(X, \mathcal{F}) = \sum\nolimits_i (-1)^i \dim_k H^i(X, \mathcal{F}).
$$
をいう。この式が意味をもつことの正当化は下を参照せよ。
\end{definition}

\noindent
定義の状況では、ベクトル空間 $H^i(X, \mathcal{F})$ のうち非零のものは
有限個しかなく（Cohomology of Spaces, Lemma
\ref{spaces-cohomology-lemma-vanishing-quasi-separated}）、その各々は有限次元である
（Cohomology of Spaces, Lemma
\ref{spaces-cohomology-lemma-proper-over-affine-cohomology-finite}）。
したがって $\chi(X, \mathcal{F}) \in \mathbf{Z}$ は良定義である。
この定義は体 $k$ に依存し、対 $(X, \mathcal{F})$ だけには依存しないことに注意せよ。

\begin{lemma}
\label{spaces-over-fields-lemma-euler-characteristic-additive}
$k$ を体とし、$X$ を $k$ 上の固有代数空間とする。
$0 \to \mathcal{F}_1 \to \mathcal{F}_2 \to \mathcal{F}_3 \to 0$ を
$X$ 上の連接加群の短完全列とする。このとき
$$
\chi(X, \mathcal{F}_2) = \chi(X, \mathcal{F}_1) + \chi(X, \mathcal{F}_3)
$$
である。
\end{lemma}

\begin{proof}
補題の短完全列に付随するコホモロジー長完全列
$$
0 \to H^0(X, \mathcal{F}_1) \to H^0(X, \mathcal{F}_2) \to
H^0(X, \mathcal{F}_3) \to H^1(X, \mathcal{F}_1) \to \ldots
$$
を考える。線型代数の階数・退化次数定理により
$$
0 = \dim H^0(X, \mathcal{F}_1) - \dim H^0(X, \mathcal{F}_2)
+ \dim H^0(X, \mathcal{F}_3) - \dim H^1(X, \mathcal{F}_1) + \ldots
$$
である。これは直ちに補題を含意する。
\end{proof}

\begin{lemma}
\label{spaces-over-fields-lemma-euler-characteristic-morphism}
$k$ を体とし、$f : Y \to X$ を $k$ 上固有な代数空間の射とする。
$\mathcal{G}$ を連接 $\mathcal{O}_Y$-加群とする。このとき
$$
\chi(Y, \mathcal{G}) = \sum (-1)^i \chi(X, R^if_*\mathcal{G})
$$
である。
\end{lemma}

\begin{proof}
この式は意味をもつ。実際、層 $R^if_*\mathcal{G}$ は連接であり、非零のものは
有限個しかない。Cohomology of Spaces, Lemmas
\ref{spaces-cohomology-lemma-proper-pushforward-coherent} および
\ref{spaces-cohomology-lemma-vanishing-higher-direct-images} を参照せよ。
Cohomology on Sites, Lemma \ref{sites-cohomology-lemma-Leray} により、
$$
E_2^{p, q} = H^p(X, R^qf_*\mathcal{G})
$$
をもち $H^{p + q}(Y, \mathcal{G})$ に収束するスペクトル系列がある。
$X$ 上のコホモロジーの有限性により、非零な $E_2^{p, q}$ は有限個しかなく、
各 $E_2^{p, q}$ は有限次元ベクトル空間である。同じことが
$E_r^{p, q}$ にも $r \geq 2$ に対して成り立ち、
$$
\sum (-1)^{p + q} \dim_k E_r^{p, q}
$$
は $r$ に依存しない。十分大きな $r$ に対して
$E_r^{p, q} = E_\infty^{p, q}$ であり、収束とは
$H^n(Y, \mathcal{G})$ 上に、その次数付き部分が
$E_\infty^{p, q}$ であり $p + 1 = n$ を満たすようなフィルトレーションが
存在することを意味する
（これがスペクトル系列の収束の意味である）ので、結論を得る。
% P09-ERR-0294
\end{proof}









\section{数値的交叉}
\label{spaces-over-fields-section-num}

\noindent
本節では、固有代数空間上の連接層の Euler 標数を用いて、可逆加群の数値的交叉数を
得るための計算を行う。主な道具は次の補題である。

\begin{lemma}
\label{spaces-over-fields-lemma-numerical-polynomial-from-euler}
$k$ を体とし、$X$ を $k$ 上の固有代数空間とする。
$\mathcal{F}$ を連接 $\mathcal{O}_X$-加群とする。
$\mathcal{L}_1, \ldots, \mathcal{L}_r$ を可逆 $\mathcal{O}_X$-加群とする。写像
$$
(n_1, \ldots, n_r) \longmapsto
\chi(X, \mathcal{F} \otimes
\mathcal{L}_1^{\otimes n_1} \otimes \ldots \otimes
\mathcal{L}_r^{\otimes n_r})
$$
は $n_1, \ldots, n_r$ の数値多項式であり、その全次数は
$\mathcal{F}$ のスキーム論的台の次元以下である。
\end{lemma}

\begin{proof}
$Z \subset X$ を $\mathcal{F}$ のスキーム論的台とする。このとき
$\mathcal{F} = i_*\mathcal{G}$ であり、ここで連接
$\mathcal{O}_Z$-加群 $\mathcal{G}$ が存在する
（Cohomology of Spaces, Lemma
\ref{spaces-cohomology-lemma-coherent-support-closed}）。また
$$
\chi(X, \mathcal{F} \otimes
\mathcal{L}_1^{\otimes n_1} \otimes \ldots \otimes
\mathcal{L}_r^{\otimes n_r}) =
\chi(Z, \mathcal{G} \otimes
i^*\mathcal{L}_1^{\otimes n_1} \otimes \ldots \otimes
i^*\mathcal{L}_r^{\otimes n_r})
$$
である。これは射影公式
（Cohomology on Sites, Lemma \ref{sites-cohomology-lemma-projection-formula}）
と Cohomology of Spaces, Lemma
\ref{spaces-cohomology-lemma-relative-affine-cohomology} による。
$|Z| = \text{Supp}(\mathcal{F})$ なので、次を示せば十分である：
$$
P_\mathcal{F}(n_1, \ldots, n_r) :
(n_1, \ldots, n_r)
\longmapsto
\chi(X, \mathcal{F} \otimes
\mathcal{L}_1^{\otimes n_1} \otimes \ldots \otimes
\mathcal{L}_r^{\otimes n_r})
$$
は $n_1, \ldots, n_r$ の数値多項式で、全次数が $\dim(X)$ 以下である。
性質 $\mathcal{P}$ が連接 $\mathcal{O}_X$-加群 $\mathcal{F}$ に対して
成り立つとは、上の主張が真であることをいう。

\medskip\noindent
この主張を dévissage により証明する。より正確には、Cohomology of Spaces, Lemma
\ref{spaces-cohomology-lemma-property-higher-rank-cohomological-variant}
の条件 (1), (2), (3) が満たされることを確認する。

\medskip\noindent
条件 (1) の確認。短完全列
$$
0 \to \mathcal{F}_1 \to \mathcal{F}_2 \to \mathcal{F}_3 \to 0
$$
をとる。これは $X$ 上の連接層の列とする。Lemma
\ref{spaces-over-fields-lemma-euler-characteristic-additive} により
$$
P_{\mathcal{F}_2}(n_1, \ldots, n_r) =
P_{\mathcal{F}_1}(n_1, \ldots, n_r) +
P_{\mathcal{F}_3}(n_1, \ldots, n_r)
$$
である。
% P09-ERR-0295
したがって層 $\mathcal{F}_i$ のうち 3 つ中 2 つが性質 $\mathcal{P}$ をもてば、
残りももつことは明らかである。

\medskip\noindent
条件 (2) は
$P_{\mathcal{F}^{\oplus m}}(n_1, \ldots, n_r) =
mP_\mathcal{F}(n_1, \ldots, n_r)$
なので従う。

\medskip\noindent
(3) の証明。$i : Z \to X$ を、$|Z|$ が既約であるような被約閉部分空間とする。
連接加群 $\mathcal{G}$ を $X$ 上に見つけなければならない。その台は $Z$ であり、
性質 $\mathcal{P}$ が $\mathcal{G}$ に対して成り立つものとする。
二つの構成を与える。一つは Chow の補題を
用い、もう一つはスキームによる有限被覆を用いる。

\medskip\noindent
スキームによる有限被覆を用いた $\mathcal{G}$ の存在証明。
% P09-ERR-0296
$\pi : Z' \to Z$ を有限全射とし、$Z'$ がスキームとなるように選ぶ。
Limits of Spaces, Proposition
\ref{spaces-limits-proposition-there-is-a-scheme-finite-over} を参照せよ。
$\mathcal{G} = i_*\pi_*\mathcal{O}_{Z'} = (i \circ \pi)_*\mathcal{O}_{Z'}$ とおく。
$Z'$ は $k$ 上固有であり、$\mathcal{G}$ の台は $Y$ であることに注意せよ
（詳細は省略する）。
% P09-ERR-0297
次が成り立つ：
$$
R(\pi \circ i)_*(\mathcal{O}_{Z'}) = \mathcal{G}
\quad\text{and}\quad
R(\pi \circ i)_*(\pi^*i^*(\mathcal{L}_1^{\otimes n_1} \otimes \ldots \otimes
\mathcal{L}_r^{\otimes n_r})
) = \mathcal{G} \otimes \mathcal{L}_1^{\otimes n_1} \otimes \ldots \otimes
\mathcal{L}_r^{\otimes n_r}
$$
% P09-ERR-0298
第 1 の等式は $i \circ \pi$ がアフィンであること
（Cohomology of Spaces, Lemma
\ref{spaces-cohomology-lemma-affine-vanishing-higher-direct-images}）から成り立ち、
第 2 の等式は第 1 の等式と射影公式
（Cohomology on Sites, Lemma \ref{sites-cohomology-lemma-projection-formula}）
から従う。Leray
（Cohomology on Sites, Lemma \ref{sites-cohomology-lemma-apply-Leray}）
を用いると
$$
P_\mathcal{G}(n_1, \ldots, n_r) =
\chi(Z', \pi^*i^*(\mathcal{L}_1^{\otimes n_1} \otimes \ldots \otimes
\mathcal{L}_r^{\otimes n_r}))
$$
を得る。スキームの場合
（Varieties, Lemma \ref{varieties-lemma-numerical-polynomial-from-euler}）により、
これは $n_1, \ldots, n_r$ の数値多項式であり、その次数は $\dim(Z')$ 以下である。
$\dim(Z') \leq \dim(Z) \leq \dim(X)$ なので結論を得る。
第 1 の不等式は Decent Spaces, Lemma
\ref{decent-spaces-lemma-dimension-quasi-finite} から従う。

\medskip\noindent
Chow の補題を用いた $\mathcal{G}$ の存在証明。
% P09-ERR-0299
Cohomology of Spaces, Lemma \ref{spaces-cohomology-lemma-weak-chow} を射
$Z \to \Spec(k)$ に適用する。すると全射固有射
$f : Y \to Z$ を得る。これは $\Spec(k)$ 上であり、$Y$ は
$\mathbf{P}^m_k$ の閉部分スキームである。ここで $m$ はある整数である。
$Y$ を閉部分スキームで置き換えれば、$Y$ は整であり、
$f : Y \to Z$ はオルタレーションであると仮定してよい。Lemma
\ref{spaces-over-fields-lemma-alteration-contained-in} を参照せよ。
$\mathcal{O}_Y(n)$ によって $\mathcal{O}_{\mathbf{P}^m_k}(n)$ の引き戻しを表す。
% P09-ERR-0300
$n > 0$ を、$R^pf_*\mathcal{O}_Y(n) = 0$ が $p > 0$ に対して成立するように選ぶ。
Cohomology of Spaces, Lemma
\ref{spaces-cohomology-lemma-kill-by-twisting} を参照せよ。
$\mathcal{G} = i_*f_*\mathcal{O}_Y(n)$ は $\mathcal{P}$ を満たすと主張する。
実際、スキームの場合
（Varieties, Lemma \ref{varieties-lemma-numerical-polynomial-from-euler}）
により
$$
(n_1, \ldots, n_r)
\longmapsto
\chi(Y, \mathcal{O}_Y(n) \otimes
f^*i^*(\mathcal{L}_1^{\otimes n_1} \otimes \ldots \otimes
\mathcal{L}_r^{\otimes n_r}))
$$
は $n_1, \ldots, n_r$ の数値多項式で、その全次数は $\dim(Y)$ 以下である。
一方、射影公式
（Cohomology on Sites, Lemma \ref{sites-cohomology-lemma-projection-formula}）により
\begin{align*}
i_*Rf_*\left(
\mathcal{O}_Y(n) \otimes
f^*i^*(\mathcal{L}_1^{\otimes n_1} \otimes \ldots \otimes
\mathcal{L}_r^{\otimes n_r})\right)
& =
i_*Rf_*\mathcal{O}_Y(n) \otimes
\mathcal{L}_1^{\otimes n_1} \otimes \ldots \otimes
\mathcal{L}_r^{\otimes n_r} \\
& =
\mathcal{G} \otimes \mathcal{L}_1^{\otimes n_1} \otimes \ldots \otimes
\mathcal{L}_r^{\otimes n_r}
\end{align*}
である。最後の等式は $n$ の選び方による。Leray
（Cohomology on Sites, Lemma \ref{sites-cohomology-lemma-apply-Leray}）
により
$$
\chi(Y, \mathcal{O}_Y(n) \otimes
f^*i^*(\mathcal{L}_1^{\otimes n_1} \otimes \ldots \otimes
\mathcal{L}_r^{\otimes n_r})) =
P_\mathcal{G}(n_1, \ldots, n_r)
$$
を得る。$\dim(Y) \leq \dim(Z) \leq \dim(X)$ なので結論を得る。
第 1 の不等式は Morphisms of Spaces, Lemma
\ref{spaces-morphisms-lemma-alteration-dimension-general} と、
$Y \to Z$ がオルタレーションであるという事実
（したがって一般点において誘導される剰余体拡大は有限である）から従う。
% P09-ERR-0301
\end{proof}

\noindent
次の補題は、最高次係数が大まかには、連接加群の台の一般点における長さのみに
依存することを示す。
% P09-ERR-0302

\begin{lemma}
\label{spaces-over-fields-lemma-numerical-polynomial-leading-term}
$k$ を体とし、$X$ を $k$ 上の固有代数空間とする。
$\mathcal{F}$ を連接 $\mathcal{O}_X$-加群とする。
$\mathcal{L}_1, \ldots, \mathcal{L}_r$ を可逆 $\mathcal{O}_X$-加群とする。
$d = \dim(\text{Supp}(\mathcal{F}))$ とおく。
$Z_i \subset X$ を $\text{Supp}(\mathcal{F})$ の次元 $d$ の既約成分とする。
$\overline{x}_i$ を $Z_i$ の幾何学的一般点とし、
$m_i = \text{length}_{\mathcal{O}_{X, \overline{x}_i}}
(\mathcal{F}_{\overline{x}_i})$ とおく。このとき
$$
\chi(X, \mathcal{F} \otimes \mathcal{L}_1^{\otimes n_1} \otimes \ldots \otimes
\mathcal{L}_r^{\otimes n_r}) -
\sum\nolimits_i
m_i\ \chi(Z_i, \mathcal{L}_1^{\otimes n_1} \otimes \ldots \otimes
\mathcal{L}_r^{\otimes n_r}|_{Z_i})
$$
は $n_1, \ldots, n_r$ の数値多項式であり、全次数は $< d$ である。
\end{lemma}

\begin{proof}
まず少し弱い主張を証明する。すなわち $\dim(X) = N$ とし、
$X_i \subset X$ を次元 $N$ の既約成分とする。
$\overline{x}_i$ を $X_i$ の幾何学的一般点とする。\'etale 局所環
$\mathcal{O}_{X, \overline{x}_i}$ は次元 $0$ の Noether 環なので、
任意の連接 $\mathcal{O}_X$-加群 $\mathcal{F}$ に対し、長さ
$$
m_i(\mathcal{F}) = \text{length}_{\mathcal{O}_{X, \overline{x}_i}}
(\mathcal{F}_{\overline{x}_i})
$$
は $\geq 0$ の整数である。次を主張する：
$$
E(\mathcal{F}) =
\chi(X, \mathcal{F} \otimes \mathcal{L}_1^{\otimes n_1} \otimes \ldots \otimes
\mathcal{L}_r^{\otimes n_r}) -
\sum\nolimits_i
m_i(\mathcal{F})\ \chi(Z_i, \mathcal{L}_1^{\otimes n_1} \otimes \ldots \otimes
\mathcal{L}_r^{\otimes n_r}|_{Z_i})
$$
% P09-ERR-0303
は $n_1, \ldots, n_r$ の数値多項式であり、全次数は $< N$ である。
これを Cohomology of Spaces, Lemma
\ref{spaces-cohomology-lemma-property-higher-rank-cohomological-variant}
を用いて証明する。任意の短完全列
$0 \to \mathcal{F}' \to \mathcal{F} \to
\mathcal{F}'' \to 0$ に対して
$E(\mathcal{F}) = E(\mathcal{F}') + E(\mathcal{F}'')$ である。
これは Euler 標数の加法性
（Lemma \ref{spaces-over-fields-lemma-euler-characteristic-additive}）と長さの加法性
（Algebra, Lemma \ref{algebra-lemma-length-additive}）から従う。
これは直ちに Cohomology of Spaces, Lemma
\ref{spaces-cohomology-lemma-property-higher-rank-cohomological-variant}
の性質 (1), (2) を含意する。最後に
$\mathcal{G} = \mathcal{O}_Z$ とすれば性質 (3) が成り立つ。
ここで $Z \subset X$ は任意の既約被約閉部分空間である。
% P09-ERR-0304
実際、$Z = Z_{i_0}$ がある $i_0$ に対して成り立つならば
$m_i(\mathcal{G}) = \delta_{i_0i}$ であり、$E(\mathcal{G}) = 0$ を得る。
$Z \not = Z_i$ がどの $i$ に対しても成り立つならば、
$m_i(\mathcal{G}) = 0$ がすべての $i$ に対して成り立ち、
$\dim(Z) < N$ なので、Lemma
\ref{spaces-over-fields-lemma-numerical-polynomial-from-euler} から結論を得る。

\medskip\noindent
補題に述べた主張の証明。
$Z \subset X$ を $\mathcal{F}$ のスキーム論的台とする。このとき
$\mathcal{F} = i_*\mathcal{G}$ であり、ここで連接
$\mathcal{O}_Z$-加群 $\mathcal{G}$ が存在する
（Cohomology of Spaces, Lemma
\ref{spaces-cohomology-lemma-coherent-support-closed}）。また
$$
\chi(X, \mathcal{F} \otimes
\mathcal{L}_1^{\otimes n_1} \otimes \ldots \otimes
\mathcal{L}_r^{\otimes n_r}) =
\chi(Z, \mathcal{G} \otimes
i^*\mathcal{L}_1^{\otimes n_1} \otimes \ldots \otimes
i^*\mathcal{L}_r^{\otimes n_r})
$$
である。これは射影公式
（Cohomology on Sites, Lemma \ref{sites-cohomology-lemma-projection-formula}）
と Cohomology of Spaces, Lemma
\ref{spaces-cohomology-lemma-relative-affine-cohomology} による。
$|Z| = \text{Supp}(\mathcal{F})$ なので、$Z_i \subset Z$ が
すべての $i$ に対して成り立ち、これらは $Z$ の次元 $d$ の既約成分である。
$\overline{x}_i$ は $Z$ の幾何点とみなすことができ、そうする。写像
$i^\sharp : \mathcal{O}_X \to i_*\mathcal{O}_Z$ は全射
$$
\mathcal{O}_{X, \overline{x}_i} \to \mathcal{O}_{Z, \overline{x}_i}
$$
を定める。この写像を介して加群の同型
$\mathcal{G}_{\overline{x}_i} = \mathcal{F}_{\overline{x}_i}$ がある。
実際 $\mathcal{F} = i_*\mathcal{G}$ である。これは
$$
m_i = \text{length}_{\mathcal{O}_{X, \overline{x}_i}}
(\mathcal{F}_{\overline{x}_i}) =
\text{length}_{\mathcal{O}_{Z, \overline{x}_i}}
(\mathcal{G}_{\overline{x}_i})
$$
を含意する。したがって補題の式は
$$
\chi(Z, \mathcal{G} \otimes \mathcal{L}_1^{\otimes n_1} \otimes \ldots \otimes
\mathcal{L}_r^{\otimes n_r}) -
\sum\nolimits_i
m_i\ \chi(Z_i, \mathcal{L}_1^{\otimes n_1} \otimes \ldots \otimes
\mathcal{L}_r^{\otimes n_r}|_{Z_i})
$$
に等しく、第 1 段落の議論を $Z$ へ、$X$ の代わりに適用すれば結論を得る。
% P09-ERR-0305
\end{proof}

\begin{definition}
\label{spaces-over-fields-definition-intersection-number}
$k$ を体とし、$X$ を $k$ 上の固有代数空間とする。
$i : Z \to X$ を次元 $d$ の閉部分空間とする。
$\mathcal{L}_1, \ldots, \mathcal{L}_d$ を可逆 $\mathcal{O}_X$-加群とする。
{\it 交点数} $(\mathcal{L}_1 \cdots \mathcal{L}_d \cdot Z)$ を
$n_1 \ldots n_d$ の係数として定義する。ここで用いる数値多項式は
$$
\chi(X, i_*\mathcal{O}_Z \otimes \mathcal{L}_1^{\otimes n_1} \otimes
\ldots \otimes \mathcal{L}_d^{\otimes n_d}) =
\chi(Z, \mathcal{L}_1^{\otimes n_1} \otimes
\ldots \otimes \mathcal{L}_d^{\otimes n_d}|_Z)
$$
である。
$\mathcal{L}_1 = \ldots = \mathcal{L}_d = \mathcal{L}$ という特別な場合には
$(\mathcal{L}^d \cdot Z)$ と書く。
\end{definition}

\noindent
定義中の表示等式は射影公式
（Cohomology, Section \ref{cohomology-section-projection-formula}）および
Cohomology of Schemes, Lemma
\ref{coherent-lemma-relative-affine-cohomology} から従う。
% P09-ERR-0306
これらの交点数についていくつかの補題を証明する。

\begin{lemma}
\label{spaces-over-fields-lemma-intersection-number-integer}
Definition \ref{spaces-over-fields-definition-intersection-number} の状況で、交点数
$(\mathcal{L}_1 \cdots \mathcal{L}_d \cdot Z)$ は整数である。
\end{lemma}

\begin{proof}
次数 $e$ の $n_1, \ldots, n_d$ に関する任意の数値多項式は、
$\mathbf{Z}$-線型結合として一意に書ける。用いる函数は
${n_1 \choose k_1}{n_2 \choose k_2} \ldots {n_d \choose k_d}$ であり、
$k_1 + \ldots + k_d \leq e$ を満たす。これを $e = d$ として適用せよ。
演習として残す。
\end{proof}

\begin{lemma}
\label{spaces-over-fields-lemma-intersection-number-additive}
Definition \ref{spaces-over-fields-definition-intersection-number} の状況で、交点数
$(\mathcal{L}_1 \cdots \mathcal{L}_d \cdot Z)$ は加法的である：
$\mathcal{L}_i = \mathcal{L}_i' \otimes \mathcal{L}_i''$ ならば
$$
(\mathcal{L}_1 \cdots \mathcal{L}_i \cdots \mathcal{L}_d \cdot Z) =
(\mathcal{L}_1 \cdots \mathcal{L}_i' \cdots \mathcal{L}_d \cdot Z) +
(\mathcal{L}_1 \cdots \mathcal{L}_i'' \cdots \mathcal{L}_d \cdot Z)
$$
である。
\end{lemma}

\begin{proof}
これは Lemma \ref{spaces-over-fields-lemma-numerical-polynomial-from-euler} により函数
\begingroup\small
$$
(n_1, \ldots, n_{i - 1}, n_i', n_i'', n_{i + 1}, \ldots, n_d)
\mapsto
\chi(Z, \mathcal{L}_1^{\otimes n_1} \otimes
\ldots \otimes (\mathcal{L}_i')^{\otimes n_i'} \otimes
(\mathcal{L}_i'')^{\otimes n_i''} \otimes \ldots \otimes
\mathcal{L}_d^{\otimes n_d}|_Z)
$$
\endgroup
が全次数 $d$ 以下の数値多項式であり、変数は $d + 1$ 個だからである。
\end{proof}

\begin{lemma}
\label{spaces-over-fields-lemma-intersection-number-in-terms-of-components}
Definition \ref{spaces-over-fields-definition-intersection-number} の状況で、
$Z_i \subset Z$ を次元 $d$ の既約成分とする。
$m_i = \text{length}_{\mathcal{O}_{X, \overline{x}_i}}
(\mathcal{O}_{Z, \overline{x}_i})$ とおく。ここで $\overline{x}_i$ は
$Z_i$ の幾何学的一般点である。このとき
$$
(\mathcal{L}_1 \cdots \mathcal{L}_d \cdot Z) =
\sum m_i(\mathcal{L}_1 \cdots \mathcal{L}_d \cdot Z_i)
$$
である。
\end{lemma}

\begin{proof}
Lemma \ref{spaces-over-fields-lemma-numerical-polynomial-leading-term} と定義から直ちに従う。
\end{proof}

\begin{lemma}
\label{spaces-over-fields-lemma-intersection-number-and-pullback}
$k$ を体とし、$f : Y \to X$ を $k$ 上固有な代数空間の射とする。
$Z \subset Y$ を次元 $d$ の整閉部分空間とし、
$\mathcal{L}_1, \ldots, \mathcal{L}_d$ を可逆 $\mathcal{O}_X$-加群とする。
このとき
$$
(f^*\mathcal{L}_1 \cdots f^*\mathcal{L}_d \cdot Z) =
\deg(f|_Z : Z \to f(Z)) (\mathcal{L}_1 \cdots \mathcal{L}_d \cdot f(Z))
$$
である。ここで $\deg(Z \to f(Z))$ は Definition \ref{spaces-over-fields-definition-degree} のとおりであり、
値を $0$ とするのは $\dim(f(Z)) < d$ の場合である。
\end{lemma}

\begin{proof}
主張における $f(Z) \subset X$ は $f$ のスキーム論的像であり、閉部分集合
$f(|Z|) \subset X$ 上の被約誘導代数空間構造でもある。Morphisms of Spaces, Lemma
\ref{spaces-morphisms-lemma-scheme-theoretic-image-reduced} を参照せよ。
% P09-ERR-0307
すると $Z$ と $f(Z)$ は $k$ 上の被約固有（したがって decent）代数空間であり、
ゆえに整である（Definition \ref{spaces-over-fields-definition-integral-algebraic-space}）。
% P09-ERR-0308
左辺は $n_1 \ldots n_d$ の係数を用いて計算される。ここで函数は
$$
\chi(Y, \mathcal{O}_Z \otimes f^*\mathcal{L}_1^{\otimes n_1} \otimes
\ldots \otimes f^*\mathcal{L}_d^{\otimes n_d}) =
\sum (-1)^i
\chi(X, R^if_*\mathcal{O}_Z \otimes
\mathcal{L}_1^{\otimes n_1} \otimes \ldots \otimes
\mathcal{L}_d^{\otimes n_d})
$$
である。
等式は Lemma \ref{spaces-over-fields-lemma-euler-characteristic-morphism} と射影公式
（Cohomology, Lemma \ref{cohomology-lemma-projection-formula}）から従う。
$f(Z)$ の次元が $< d$ ならば、Lemma
\ref{spaces-over-fields-lemma-numerical-polynomial-from-euler} により右辺は全次数 $< d$ の
多項式であり、結論は真である。$\dim(f(Z)) = d$ と仮定する。
すると次元論（Lemma \ref{spaces-over-fields-lemma-alteration-dimension}）により、Lemma
\ref{spaces-over-fields-lemma-finite-degree} の同値な条件 (1) -- (5) が成り立つ。
したがって $\deg(Z \to f(Z))$ は良定義である。
すでに用いた Lemma \ref{spaces-over-fields-lemma-finite-degree} により、
$f : Z \to f(Z)$ は空でない開部分集合 $V$ 上で有限であり、
その開部分集合は $f(Z)$ に含まれる。
$V$ を必要なら縮小し、$V$ はスキームであると仮定してよい。
$\xi \in V$ を一般点とする。このとき
$\deg(f : Z \to f(Z))$ は $f_*\mathcal{O}_Z$ の $\xi$ における茎の
$\mathcal{O}_{X, \xi}$ 上の長さであり、
% P09-ERR-0309
% P09-ERR-0310
$R^if_*\mathcal{O}_X$ の $\xi$ における茎は $i > 0$ に対して零である
% P09-ERR-0311
（たとえば Cohomology of Spaces, Lemma
\ref{spaces-cohomology-lemma-finite-higher-direct-image-zero} による）。
したがって各項
$\chi(X, R^if_*\mathcal{O}_Z \otimes
\mathcal{L}_1^{\otimes n_1} \otimes \ldots \otimes
\mathcal{L}_d^{\otimes n_d})$ は $i > 0$ に対して全次数 $< d$ であり、
$$
\chi(X, f_*\mathcal{O}_Z \otimes
\mathcal{L}_1^{\otimes n_1} \otimes \ldots \otimes
\mathcal{L}_d^{\otimes n_d})
=
\deg(f : Z \to f(Z)) \chi(f(Z),
\mathcal{L}_1^{\otimes n_1} \otimes \ldots \otimes
\mathcal{L}_d^{\otimes n_d}|_{f(Z)})
$$
が、全次数 $< d$ の多項式を法として成り立つ。これは Lemma
\ref{spaces-over-fields-lemma-numerical-polynomial-leading-term} による。求める結論が従う。
\end{proof}

\begin{lemma}
\label{spaces-over-fields-lemma-numerical-intersection-effective-Cartier-divisor}
$k$ を体とし、$X$ を $k$ 上の固有代数空間とする。
$Z \subset X$ を次元 $d$ の閉部分空間とする。
$\mathcal{L}_1, \ldots, \mathcal{L}_d$ を
可逆 $\mathcal{O}_X$-加群とする。ある有効 Cartier 因子
$D \subset Z$ が存在して
$\mathcal{L}_1|_Z \cong \mathcal{O}_Z(D)$ を満たすと仮定する。このとき
$$
(\mathcal{L}_1 \cdots \mathcal{L}_d \cdot Z) =
(\mathcal{L}_2 \cdots \mathcal{L}_d \cdot D)
$$
である。
\end{lemma}

\begin{proof}
$X$ を $Z$ で、$\mathcal{L}_i$ を $\mathcal{L}_i|_Z$ で置き換えてよい。
したがって $X = Z$ および $\mathcal{L}_1 = \mathcal{O}_X(D)$ と仮定してよい。
このとき $\mathcal{L}_1^{-1}$ は $D$ のイデアル層であり、短完全列
$$
0 \to \mathcal{L}_1^{\otimes -1} \to \mathcal{O}_X \to \mathcal{O}_D \to 0
$$
を考えることができる。
$$
P(n_1, \ldots, n_d) =
\chi(X, \mathcal{L}_1^{\otimes n_1} \otimes \ldots \otimes
\mathcal{L}_d^{\otimes n_d})
$$
および
$$
Q(n_1, \ldots, n_d) =
\chi(D, \mathcal{L}_1^{\otimes n_1} \otimes \ldots \otimes
\mathcal{L}_d^{\otimes n_d}|_D)
$$
とおく。加法性
（Lemma \ref{spaces-over-fields-lemma-euler-characteristic-additive}）から
$$
P(n_1, \ldots, n_d) - P(n_1 - 1, n_2, \ldots, n_d) =
Q(n_1, \ldots, n_d)
$$
を得る。$P$ の全次数は高々 $d$ なので、
$n_1 \ldots n_d$ の $P$ における係数は
$n_2 \ldots n_d$ の $Q$ における係数に等しい。
\end{proof}
